%9Summer 22/7 Cut factorizations
\documentclass{article}
%\documentclass[hypertex,11pt]{article}
%\usepackage{showlabels}
\usepackage{amsfonts,amsmath,amscd,amssymb,amsthm}
\usepackage{stackrel}
\usepackage{fullpage}
\usepackage{xcolor}
\usepackage[T1]{fontenc}
\usepackage{imakeidx}
\usepackage{subcaption}
\makeindex[columns=3 title=Alphabetical Index]
\makeindex
%\usepackage{showlabels}
\usepackage[T2A,T1]{fontenc}
\usepackage[utf8]{inputenc}
%\usepackage[pagebackref]{hyperref}
%\usepackage[backrefs, alphabetic,initials]{amsrefs}%use this with \input{shorthead-ams-gsm.tex}
%\usepackage[alphabetic,initials]{amsrefs}%this kills the up-arrows in the bibliography 
\usepackage[all]{xy}
%\usepackage{hyperref} %prob

\usepackage[alphabetic,initials]{amsrefs}%\usepackage{mathrsfs}
\input xy%
\usepackage{fullpage}
\usepackage{amsrefs}
\makeindex
 \def\emptyline{\vspace{12pt}}
%\usepackage{fullpage}
\textheight 21.6cm \textwidth 14cm
\def\bydefn{\stackrel{def}{=}}\def\convf{\hbox{\space \raise-2mm\hbox{$\textstyle      \bigotimes \atop \scriptstyle \omega$} \space}}
\def\tl{{\tilde \gl}} 
\def\0{{\bar 0}}
\def\tm{{\tilde \mu}} 
\def\1{{\bar 1}}
\def\a{{\mathtt a}}
\def\b{{\mathtt b}}
\def\c{{\mathtt c}}
\def\z{{\mathtt z}}
\def\q{{\mathtt q}}
\def\C{{\mathbb C}}
\def\Z{{\mathbb Z}}
\def\Q{{\mathbb Q}}
\def\B{{\mathbb B}}
\def\vv{{\mathbb V}}
\def\vw{{\mathbb W}}
\def\E{{\mathbb E}}
\def\T{{\mathbb T}}
\def\G{{\mathbb G}}
%\newcommand{\Z}{{\mathbb{Z}}
\def\u{\overline{u}}
\def\w{\overline{w}}
\def\ww{\wedge}
\def\N{{\mathbb N}}
%\def\P{{\mathcal P}}
%\def\P{{\bf P}}
\def\P{{\mathbb P}}
\def\H{{\mathscr H}}
\def\R{{\mathbb R}}
\def\I{{\mathbb I}}%\def\B{{\mathscr B}}
\def\V{\mathtt{V}}
\def\xyz{{X\times_{Y}Y'}}
\def\X{{\mathbb X}}
\def\U{\mathscr{U}}
\def\W{\mathscr{W}}
\def\D{\mathbb{D}}
%\def\G{\mathscr{G}}
\def\degf{{\rm\operatorname{deg_{filt}}}}
\def\Fract {{\operatorname{Fract}}}
\def\pl{{\operatorname{pole}}}
\def\zero{{\operatorname{zero}}}
\def\Max{{\operatorname{Max}\;}}
\def\noniso{{\operatorname{non-isotropic}}}
\def\GrSpec{{\operatorname{GrSpec}\;}}
\def\GrPrim{{\operatorname{GrPrim}\;}}
\def\Ber{{\operatorname{Ber}\;}}
\def\Div{{\operatorname{Div_\Z}}}
\def\int{{\operatorname{Int_\Z}}}
%\def\Stab{{\operatorname{Stab}\;}}
%\def\Cent{{\operatorname{Cent}\;}}
\def\du{\stackrel{\cdot}{\bigcup}}
\def\ds{\stackrel{\cdot}{\cup}}
\def\Lie{{\operatorname{Lie}}}
\def\diag{{\operatorname{diag}}}
\def\height{{\operatorname{ht}}}
\def\odd{{\operatorname{odd}}}
\def\Stab{{\operatorname{Stab}}}
\def\Cent{{\operatorname{Cent}}}
\def\tors{{\operatorname{tors}}}
\def\Qb{{\operatorname{\frak{Qco}\; \X_\gb}}}
\def\Qc{{\operatorname{\frak{Qco}\; \X_\gc}}}
\def\Qbb{{\operatorname{\frak{Qco}\; \X_{\gb,\gb}}}}
\def\Qcc{{\operatorname{\frak{Qco}\; \X_{\gc,\gc}}}}
\def\Qbc{{\operatorname{\frak{Qco}\; \X_{\gb,\gc}}}}
\def\Qcb{{\operatorname{\frak{Qco}\; \X_{\gc,\gb}}}}


\def\Obb{{\operatorname{\cO_{\gb,\gb}}}}
\def\Occ{{\operatorname{\cO_{\gc,\gc}}}}
\def\Obc{{\operatorname{\cO_{\gb,\gc}}}}
\def\Ocb{{\operatorname{\cO_{\gc,\gb}}}}
\def\Spec{{\operatorname{Spec}\;}}
\def\Prim{{\operatorname{Prim}\;}}
\def\odet{{\operatorname{odet \;}}}
\def\otr{{\operatorname{otr \;}}}
\def\ad{{\operatorname{ad \;}}}
\def\ado{{\operatorname{\overline{ad} \;}}}
\def\Ad{{\operatorname{Ad \;}}}
\def\trp{{\operatorname{transpose}}}
\def\mo{{\operatorname{Mor}}}
\def\hgt{{\operatorname{ht}}}
\def\top{{\operatorname{top}}}
\def\inr{{\operatorname{ind}}}
\def\rer{{\operatorname{res}}}
\def\ind{{\operatorname{ind}_^G\;}}
\def\rad{\operatorname{rad\;}}
\newcommand{\lleft}{\longleftarrow}
\def\gra{{\operatorname{gr}\;}}
\def\gra{{\operatorname{gr}\;}}
\def\Trace{{\operatorname{Trace}\;}}
\def\ST{{\operatorname{Str}\;}}
\def\Res{{\operatorname{Res}}}
\def\Ind{{\operatorname{Ind}}}
\def\df{{\operatorname{def}\;}}
\def\cdf{{\operatorname{codef}}}
\def\can{{\operatorname{can}}}
\def\Int{{\operatorname{Int}}}
\def\lh{{\operatorname{\log_h\;}}}
\def\res{{\operatorname{res}_P^G\;}}

\def\hy{\operatorname{-mod}}

\def\Der {{\operatorname{Der}\;}}
\def\Aut {{\operatorname{Aut}\;}}
\def\Out {{\operatorname{Out}\;}}
\def\ch{{\operatorname{ch}\:}}
\def\wtsp{{\operatorname{wt}\;}}
\def\pwg1{{\operatorname{PWG}}}
\def\pwg{{\operatorname{pwg}}}
\def\pro{{\operatorname{Proj\,}}}
\def\Pr{{\operatorname{\bf Proj}\;}}
\def\gldim{{\operatorname{gl \; dim}\;}}
\def\injdim{{\operatorname{inj \; dim}\;}}
\def\prd{{\operatorname{proj \; dim}\;}}
\def\idd{{\operatorname{ind_P^G}\;}}
\def\ind{{\operatorname{ind}_P^G\;}}
\def\res{{\operatorname{res}_P^G\;}}
\def\trace{{\operatorname{tr}\;}}
\def\resr{{\operatorname{res}\;}}
\def\der {{\operatorname{Der}}}
\def\Spin{{\operatorname{Spin}}}
\def\cf {{\operatorname{cf}}}
\def\indr{{\operatorname{ind}\;}}
\def\coind{{\operatorname{coind}\;}}
\def\wt{{\operatorname{wt}}}
\def\rank{{\operatorname{rank}}}
\def\span{{\operatorname{span}}}
\def\Tor{{\operatorname{Tor}}}
\def\Re{{\operatorname{Re}}}
\def\CS{{\operatorname{C-span}}\;}
\def\tr {{\operatorname{transpose}}}
\newcommand{\ttx}{\mathtt{x}}
\newcommand{\ttv}{\mathtt{v}}
%\def\ht{{\operatorname{ht}}}
\def\Ext{{\operatorname{Ext}}}
\def\Homv {{\operatorname{Hom}}}%A \Hom variant use for Hom_R
\def\Hom {{\operatorname{Hom}}}%use for Hom (ie Hom_K)
\def\Ker {{\operatorname{Ker}\;}}
\def\End {{\operatorname{End}}}
\def\Mod {{\operatorname{Mod-}}}
\def\id {{\operatorname{id}}}
\def\ev {{\operatorname{ev}}}
\def\Gmod {G{\operatorname{-{\bf mod}}}}
\def\Gomod {(\fg, G_0){\operatorname{-{\bf mod}}}}
\def\triv {{\operatorname{triv}}}
\def\dist{{\operatorname{dist}}}
\def\cl{{\operatorname{close}}}
\def\ann {{\operatorname{ann}}}
\def\rann {{\operatorname{r.ann}}}
\def\core {{\operatorname{core}}}
\def\lann {{\operatorname{lann}}}
\def\Sder {{\operatorname{Sder}}}
\def\Endv {{\operatorname{End}}}%A \End variant use for End_R
%\def\rank{{\operatorname{rank}\;}}
\def\Comm{{\operatorname{Comm}\;}}
\def\Img{{\operatorname{Im}\;}}
\def\arctanh{{\operatorname{arctanh}\;}}
\newcommand{\tth}{\mathtt{hy}}
\newcommand{\ttk}{\mathtt{k}}
\newcommand{\tte}{\mathtt{e}}
\newcommand{\ttd}{\mathtt{d}}
\newcommand{\ttS}{\mathtt{S}}
\newcommand{\ttT}{\mathtt{T}}
\newcommand{\wo}{{\stackrel{o}{W}}}
\newcommand{\bk}{{\bf k}}
\newcommand{\lla}{L(\gl)}
\newcommand{\mla}{M(\Lambda)}
%\newcommand{\ovd}{\overline{\Delta}}
\newcommand{\ovd}{{\Delta}}
\newcommand{\ws}{W^{\#}}
\newcommand{\gL}{\Lambda}
\newcommand{\gC}{\Gamma}
\newcommand{\fgm}{\stackrel{{\rm _o}}{\fg}^{-ZZ}}
\newcommand{\fgtm}{\stackrel{{\rm o}}{\;\;\fg^{-\gth}}}
\newcommand{\fgo}{\stackrel{{\rm _o}}{\fg}}
\newcommand{\pigo}{\stackrel{{\rm _o}}{\Pi}}
\newcommand{\pigov}{\stackrel{{\rm _o}}{\Pi^\vee}}
\newcommand{\pigop}{\stackrel{{\rm _o}}{\Pi'}}
\newcommand{\fga}{\stackrel{{\rm _o}}{\fg}^\ga}
\newcommand{\fgt}{\stackrel{{\rm _o}}{\fg}^\gth}
\newcommand{\fgaa}{\stackrel{{\rm o}}{\;\;\fg\;^\ga}}
\newcommand{\Dto}{\stackrel{{\rm _o}}{\Delta}}
\newcommand{\fho}{\stackrel{{\rm _o}}{\fh}}

\newcommand{\fgk}{\underline{\fg}_k}
%\newcommand{\fgk}{\stackrel{{\rm _k}}{\underline{\fg}}}
%\newcommand{\fhs}{\stackrel{{\rm _o}}{\fh}^*}
\newcommand{\efl}{e_0, e_1, \ldots, e_n, f_0, f_1, \ldots, f_n}
\newcommand{\Lgo}{L(\stackrel{{\rm _o}}{\fg})}
\newcommand{\Lgh}{\widehat{L}(\stackrel{{\rm _o}}{\fg})}
\newcommand{\itemi}{\item[{{\rm(i)}}]}
\newcommand{\itemii}{\item[{{\rm(ii)}}]}
\newcommand{\itemiii}{\item[{{\rm(iii)}}]}
\newcommand{\itemiv}{\item[{{\rm(iv)}}]}
\newcommand{\itemv}{\item[{{\rm(v)}}]}
\newcommand{\itemvi}{\item[{{\rm(vi)}}]}
\newcommand{\itema}{\item[{{\rm(a)}}]}
\newcommand{\itemb}{\item[{{\rm(b)}}]}
\newcommand{\itemc}{\item[{{\rm(c)}}]}
\newcommand{\itemd}{\item[{{\rm(d)}}]}
\newcommand{\iteme}{\item[{{\rm(e)}}]}
\newcommand{\itemf}{\item[{{\rm(f)}}]}
\newcommand{\itemg}{\item[{{\rm(g)}}]}
\newcommand{\itemh}{\item[{{\rm(h)}}]}
\newcommand{\da}{\downarrow}
\newcommand{\itemo}{\item[{}]}
\newcommand{\noi}{\noindent}
\newcommand{\ga}{\alpha}
\newcommand{\gb}{\beta}
\newcommand{\gc}{\gamma}
\newcommand{\gcc}{\overline{\gamma}}
\newcommand{\Gc}{\Gamma}
\newcommand{\Go}{\Omega}
\newcommand{\Gd}{\Delta}
\newcommand{\Gt}{\Theta}
\newcommand{\gd}{\delta}
\newcommand{\gph}{\phi}
\newcommand{\eps}{\epsilon}
\newcommand{\gs}{\sigma}
\newcommand{\gS}{\Sigma}
\newcommand{\go}{\omega}
\newcommand{\gO}{\Omega}
\newcommand{\cg}{{\rm Cdeg}}
\newcommand{\gt}{\tau}
\newcommand{\gz}{\zeta}
\newcommand{\gi}{\iota}
\newcommand{\gl}{\lambda}
\newcommand{\la}{\lambda}
\newcommand{\La}{\Lambda}
\newcommand{\gr}{\rho}
\newcommand{\gk}{\kappa}
\newcommand{\gep}{\epsilon}
\newcommand{\gth}{\theta}
\newcommand{\op}{\oplus}
\newcommand{\vph}{\varphi}
\newcommand{\gve}{\varepsilon}
\newcommand{\A}{\mathbb A}
\def\pr{{\operatorname{pr}}}
\def\Shff{{\operatorname{Shff}}}
\def\LT{\operatorname{LT}}
\def\Supp{{\operatorname{Supp}}}
\def\Im{{\operatorname{Im}\;}}
\def\LL{{\mathtt{L}}}
\def\RR{{\mathtt{R}}}
\def\CC{{\mathtt{C}}}
\newcommand{\ot}{\otimes}


\newcommand{\spp}{\mathfrak{sp}}
\newcommand{\fg}{\mathfrak{g}}\newcommand{\fgl}{\mathfrak{gl}}\newcommand{\so}{\mathfrak{so}}
\newcommand{\fsl}{\mathfrak{sl}}\newcommand{\fpsl}{\mathfrak{psl}}\newcommand{\osp}{\mathfrak{osp}}\newcommand{\spo}{\mathfrak{spo}}
\newcommand{\fgf}{\mathfrak{f}}
\newcommand{\fr}{\mathfrak{r}}
\newcommand{\fu}{\mathfrak{u}}
\newcommand{\fv}{\mathfrak{v}}
\newcommand{\fw}{\mathfrak{w}}
\newcommand{\fs}{\mathfrak{s}}
\newcommand{\ft}{\mathfrak{t}}
\newcommand{\fh}{\mathfrak{h}}
\newcommand{\fz}{\mathfrak{z}}
\newcommand{\fb}{\mathfrak{b}}
\newcommand{\fm}{\mathfrak{m}}
\newcommand{\fn}{\mathfrak{n}}
\newcommand{\fa}{\mathfrak{a}}
\newcommand{\fc}{\mathfrak{c}}
\newcommand{\fo}{\mathfrak{o}}
\newcommand{\fsp}{\mathfrak{sp}}
\newcommand{\fk}{\mathfrak{k}}
\newcommand{\fp}{\mathfrak{p}}
\newcommand{\fq}{\mathfrak{q}}
\newcommand{\fl}{\mathfrak{l}}
\newcommand{\sos}{\small{\mathfrak{osp}}}
\newcommand{\sfh}{\small{\mathfrak{h}}}
\newcommand{\sfp}{\small{\mathfrak{p}}}
\newcommand{\sfq}{\small{\mathfrak{q}}}
\newcommand{\sfb}{\small{\mathfrak{b}}}
\newcommand{\sfr}{\small{\mathfrak{r}}}
\newcommand{\sfg}{\small{\mathfrak{g}}}
\newcommand{\sfl}{\small{\mathfrak{l}}}
\newcommand{\sfn}{\small{\mathfrak{n}}}
\newcommand{\sfk}{\small{\mathfrak{k}}}
\newcommand{\sfs}{\small{\mathfrak{s}}}
\newcommand{\sfgl}{\small{\mathfrak{gl}}}
\newcommand{\ff}{\footnote}
\newfont{\eufm}{eufm10 scaled\magstep1}
%\newcommand{\mathfrak}[1]{\mbox{\eufm#1}}
\newcommand{\pd}{\partial}
\newcommand{\ci}{\circ} \newcommand{\ti}{\times}
\newcommand{\bv}{\bigvee}
\newcommand{\bw}{\bigwedge}
\newcommand{\wg}{\wedge}
\newcommand{\cO}{\mathcal{O}}
\newcommand{\cC}{\mathcal{C}}
\newcommand{\cd}{\mathcal{d}}
\newcommand{\cI}{\mathcal{I}}
\newcommand{\cN}{\mathcal{N}}
\newcommand{\cD}{\mathcal{D}}
\newcommand{\cU}{\mathcal{U}}
\newcommand{\cB}{\mathcal{B}}
\newcommand{\cjx}{\mathcal{J}_X}
\newcommand{\cA}{\mathcal{A}}
\newcommand{\cM}{\mathcal{M}}
\newcommand{\cP}{\mathcal{P}}
\newcommand{\cF}{\mathcal{F}}
\newcommand{\cY}{\mathcal{Y}}
\newcommand{\cE}{\mathcal{E}}
\newcommand{\cL}{\mathcal{L}}
\newcommand{\cG}{\mathcal{G}}
\newcommand{\cK}{\mathcal{K}}
\newcommand{\ii}{\mathcal{I}}
\newcommand{\cS}{\mathcal{S}}
\newcommand{\cT}{\mathcal{T}}
\newcommand{\cR}{\mathcal{R}}
\newcommand{\cQ}{\mathcal{Q}}
\newcommand{\CQ}{\overline{\cQ}}
%\newcommand{\cY}{\mathcal{Y}}
\newcommand{\cV}{\mathcal{V}}
\newcommand{\cH}{\mathcal{H}}
\newcommand{\cJ}{\mathcal{J}}
\newcommand{\cW}{\mathcal{W}}
\newcommand{\bbN}{\mathbb{N}}
\newcommand{\bbZ}{\mathbb{Z}}
\newcommand{\bbC}{\mathbb{C}}
\newcommand{\bbR}{\mathbb{R}}
\newcommand{\bbQ}{\mathbb{Q}}
\newcommand{\bbI}{\mathbb{I}}
\newcommand{\bbA}{\mathbb{A}}
\newcommand{\bbP}{\mathbb{P}}
\newcommand{\bbF}{\mathbb{F}}
\newcommand{\bbG}{\mathbb{G}}
\newcommand{\bbB}{\mathbb{B}}
\newcommand{\bbX}{\mathbb{X}}
\newcommand{\bbY}{\mathbb{Y}}
\newcommand{\bbJ}{\mathbb{I}}
\def\K{{\mathscr K}}
\newcommand{\ey}{\end{eqnarray}}
\newcommand{\by}{\begin{eqnarray}}
\newcommand{\nn}{\nonumber}
\newcommand{\bs}{Borel subalgebra \;}
\newcommand{\bss}{Borel subalgebras \;}
\newcommand{\bco}{\begin{conjecture}}
\newcommand{\ba}{\begin{alg}}
\newcommand{\ea}{\end{alg}}
\newcommand{\eco}{\end{conjecture}}
\newcommand{\bpf}{\begin{proof}}
\newcommand{\epf}{\end{proof}}
\newcommand{\bt}{\begin{theorem}}
\newcommand{\abt}{\begin{theorem a}}
\newcommand{\aet}{\end{theorem a}}
\newcommand{\et}{\end{theorem}}
\newcommand{\csa}{Cartan subalgebra \;}
\newcommand{\br}{\begin{rem}}
\newcommand{\er}{\end{rem}}
\newcommand{\brs}{\begin{rems}}
\newcommand{\ers}{\end{rems}}
\newcommand{\bi}{\begin{itemize}}
\newcommand{\ei}{\end{itemize}}
\newcommand{\bl}{\begin{lemma}}
\newcommand{\bsul}{\begin{sublemma}}
\newcommand{\esul}{\end{sublemma}}
\newcommand{\bp}{\begin{proposition}}
\newcommand{\be}{\begin{equation}}
\newcommand{\bc}{\begin{corollary}}
\newcommand{\bexs}{\begin{examples}}
\newcommand{\eexs}{\end{examples}}
\newcommand{\bexa}{\begin{example}}
\newcommand{\eexa}{\end{example}}
\newcommand{\bex}{\begin{exercise}}
\newcommand{\eex}{\end{exercise}}
\newcommand{\btab}{\begin{tab}}
\newcommand{\etab}{\end{tab}}
\newcommand{\el}{\end{lemma}}
\newcommand{\ep}{\end{proposition}}
\newcommand{\ee}{\end{equation}}
\newcommand{\ec}{\end{corollary}}
\newcommand{\Bc}{\begin{center}}
\newcommand{\Ec}{\end{center}}
\newcommand{\bh}{\begin{hyp}}
\newcommand{\eh}{\end{hyp}}
\newcommand{\bhs}{\begin{hyps}}
\newcommand{\ehs}{\end{hyps}}
\newcommand{\bd}{\begin{dfn}}
\newcommand{\ed}{\end{dfn}}
\newcommand{\se}{\subseteq}
\newcommand{\bca}{\bigcap}
\newcommand{\bcu}{\bigcup}
\newcommand{\bop}{\bigoplus}
\newcommand{\bsk}{\backslash}

%\newcommand{\ey}{\end{eqnarray}}
%\newcommand{\by}{\begin{eqnarray}}
%\newcommand{\el}{\end{lemma}}
%\newcommand{\ep}{\end{proposition}}
%\newcommand{\ee}{\end{equation}}
%\newcommand{\ec}{\end{corollary}}
\begin{document}
%\title{Table of Contents}
\newtheorem{thm}{Theorem}[section]
%\newtheorem{conjecture}{thm}[conjecture]
%\newtheorem{lem}[thm]{Lemma}
\newtheorem{hyp}[thm]{Hypothesis}
 \newtheorem{hyps}[thm]{Hypotheses}
%  \newtheorem{hyp}[theorem]{Hypothesis}
  %\newtheorem{rem}[theorem]{Remark}
  \newtheorem{rems}[thm]{Remarks}

\newtheorem{conjecture}[thm]{Conjecture}
\newtheorem{theorem}[thm]{Theorem}
\newtheorem{theorem a}[thm]{Theorem A}
\newtheorem{example}[thm]{Example}
\newtheorem{examples}[thm]{Examples}
\newtheorem{corollary}[thm]{Corollary}
%\newtheorem{def}[thm]{Definition}
\newtheorem{rem}[thm]{Remark}
\newtheorem{lemma}[thm]{Lemma}
\newtheorem{sublemma}[thm]{Sublemma}
\newtheorem{cor}[thm]{Corollary}
\newtheorem{proposition}[thm]{Proposition}
%\newtheorem{defi}[thm]{Definition}
%\newtheorem{rem}[thm]{Remark}
\newtheorem{exs}[thm]{Examples}
\newtheorem{ex}[thm]{Example}
%\newtheorem{al}[thm]{Algorithm}
\newtheorem{exercise}[thm]{Exercise}
%\newtheorem{part}[thm]{Exercise}
\numberwithin{equation}{section}%
\setcounter{part}{0}
%\numberwithin{exercise}{section}
%\setcounter{part}{0}
\newcommand{\drar}{\rightarrow}
\newcommand{\lra}{\longrightarrow}
\newcommand{\rra}{\longleftarrow}
\newcommand{\dra}{\Rightarrow}
\newcommand{\dla}{\Leftarrow}
\newcommand{\xra}{\xrightarrow{\sim}}
\def\sd{\operatorname{SDer}}
\def\Str{{\operatorname{Str}\;}}
%\def\drar{\raisebox{1pt}{\small\mbox{$\ssc-\cdots\rar$}}}

\newtheorem{Thm}{Main Theorem}



%-----------------------------------------------------------------

%\newtheorem{thm}{Theorem}[section]
\newtheorem*{thm*}{Theorem}
\newtheorem{lem}[thm]{Lemma}
\newtheorem*{lem*}{Lemma}
%\newtheorem{prop}[thm]{Proposition}
\newtheorem*{prop*}{Proposition}
%\newtheorem{cor}[thm]{Corollary}
\newtheorem*{cor*}{Corollary}
\newtheorem{dfn}[thm]{Definition}
\newtheorem*{defn*}{Definition}
\newtheorem{notadefn}[thm]{Notation and Definition}
\newtheorem*{notadefn*}{Notation and Definition}
\newtheorem{nota}[thm]{Notation}
\newtheorem*{nota*}{Notation}
\newtheorem{note}[thm]{Remark}
%\newtheorem{dfn}[thm]{Definition}
\newtheorem*{note*}{Remark}
\newtheorem*{notes*}{Remarks}
\newtheorem{hypo}[thm]{Hypothesis}
%\newtheorem{ex}[thm]{Example}
\newtheorem*{ex*}{Example}
\newtheorem{prob}[thm]{Problems}
\newtheorem{conj}[thm]{Conjecture}
%\begin{document}
%
%\setcounter{page}{0}

%\end{document}
\title{Geometry and coefficients of Šapovalov elements for KM  Lie superalgebras}
\author{Ian M. Musson\\%Department of Mathematical Sciences\\
University of Wisconsin-Milwaukee\\ email: {\tt
musson@uwm.edu}}


\maketitle 

\begin{abstract} 
We study \v Sapovalov   elements for symmetrizable, integrable Kac-Moody 
algebras $\fg$. 
Let  $\gc$ be a positive  root of $\fg$ and $m$ a positive integer, with suitable conditions on the pair   $(\gc, m).$ The \v Sapovalov element
$\theta_{\gamma,m}\in U(\fb^-)	$
has the important property that if $\gl$ lies on a certain hyperplane, then
$\theta_{\gc, m} v_\lambda$ is a highest weight vector of weight $\lambda -m\gc$ in
$M(\gl)$. We give  a closed fomula for all coefficients that arise in the induction step of the 
proof. The commutative version  of this formula relates the leading terms  using maps graded  algebras that we call of \v Sapovalov   algebras. This suggests  a  geometric setting  for the study of 
\v Sapovalov elements. %For Lie superalgebras comparable results hold for real roots. %under some mild restrictions.
\end{abstract}
\noi \ref{uv.1}  {Introduction}%2\ref{PA.18}
\\ 
 \noi  \ref{sss7.1}  {Preliminaries} \\
\noi{\ref{PA.2} {Main results on coefficients}\\
\noi \ref{1s.5}} {Proofs of Theorems   \ref{1sag} and  \ref{cot}}\\
\noi \ref{g1} {Commutative version}\\ 
\noi  \ref{GF} {Geometry of Šapovalov elements} \\ 
\noi \ref{vf} Examples

%\end{document}


\section{Introduction.} \label{uv.1} %\ref{Main Results.}.

Using  a remarkable combination of representation theory and algebraic geometry,   \v Sapovalov   elements and the \v Sapovalov   determinant  were introduced in \cite{Sh}.  A few examples suffice to demonstrate their importance. 
The \v Sapovalov determinant is used to construct the Jantzen filtration for semisimple Lie algebras \cite{J1} and also to compute the center of the enveloping algebra of a finite dimensional Kac-Moody Lie superalgebra \cite{Gk}, \cite{Kac4}.
 \v Sapovalov elements are crucial for understanding maps between Verma modules. Fix a positive root $\gb$ and integer $m$.  Define a \index{Hyperplane ${H}_{\gb, m}$ }hyperplane 
\be \label{vat}{H}_{\gb, m} = \{ \lambda \in  {\mathfrak h}^*|(\lambda+\gr, \gb) = m(\gb, \gb)/2  \}.\ee  %In this paper the Verma module $M(\gl)$ is generated by a highest weight vector $v_\gl$ of weoght $\gl-\gt$. 
%Except in Subsection \ref{Type C}. t
The Verma module $M(\gl)$ is generated by a highest weight vector $v_\gl$ of weight $\gl$.
\index{Verma module $M(\gl)$ with highest  weight $\gl$} 
We postpone some of the definitions in this introduction and for simplicity assume $\fg$ is a simple Lie algebra. A {\it \v Sapovalov element} $\theta_{\gb, m} \in U(\fb^-)$  has the property that  for all
$\gl \in H_{\gb, m}$, 
$\theta_{\gb, m} v_\lambda$ is a highest weight vector of weight $\lambda -m\gb=\mu$ in
$M(\gl)$.%\ff{Not true for superalgebras, for ex $\fsl(2,1)$ when $\gl =  -\gr =\gb.$ Without this condition there is no map} IAN
\\ \\
Define a 
{\it \v Sapovalov map} between Verma modules $M(\mu)\lra M(\gl)$  by $xv_\mu \lra x\theta_{\gb, m} v_\gl $.  In the  Lie algebra case any nonzero map between Vermas is a composite of \v Sapovalov maps. If $\gb_0$ is a simple root, then $\theta_{\gb_0,  m} = e_{-\gb_0}^m$. %In general to construct $\theta_{\gb,m}$, we write $\gb = u\gb_0$ with $u$ a Weyl group element of shortest length. 
Thus what is of interest in the construction of \v Sapovalov  elements
is the induction step.  Suppose that 
\be \label{vct}
\gc = s_\ga \gb = \gb +q\ga\ee for some $q>0.$ If   $\theta_{\gb,m}$ has been constructed we need to find 
$\theta_{\gc,m}$.  %This involves localization at the negative root vector $e=e_{-\ga}$. %associated to $s_\ga$,
\\ \\
Many of the results of this paper come in commutative and noncommutative  versions.  Let us explain what this means.  In Theorem \ref{cgw} using  two linear maps $B_U$ and $\nabla^{t}_U$ %\ff{At first $t$ is the integer that  appears in \eqref{121d}, but by Zariski density we can then use any $t\in \ttk$}, 
we give a closed formula 
for all the coefficients 
of  $\theta_{\gc,m}$ in terms of the coefficients of 
$\theta_{\gb,m}$. %Note that  $\nabla_U$ is a map $\nabla_U: U(\fn^-)_e\lra  U(\fn^-)_e[T] $. By evaluation at $T=t$,  where  $t$ is as in \eqref{121d} we obtain a map  $\nabla^{t}_U: U(\fn^-)_e \lra U(\fn^-)_e$. 
Theorem \ref{cgw} depends on a particular ordering on the set of positive roots, and thus it makes sense to pass to the commutative associated graded ring.  
\subsubsection{The symbol  map and symmetrization} \label{syb}
First a clarification. The  symmetrization map 
$\go:S(\fn^-)\lra U(\fn^-)$ is defined in Subsection \ref{ss6.4}.
For nonzero $a\in S(\fh)$, write $a=\sum_{i=0}^n a_i$ where $a_i$ is homogeneous of degree $i$ and  $a_n\neq 0.$  Then define the {\it symbol}  \index{Symbol map $\gk$} of $a$ to be $\gk(a)= a_n \in S(\fh)$.  Thus if $n>0$,   $\deg (\gk(a)-a)<n.$ 
We use a basis 
$\{x^\ttk_{\gs}\}$ of  $S(\fn^-)$ indexed by partitions, see
\eqref{kpar}. Then any $x\in U(\fb^-)$ can be written uniquely as
\be \label{nfn}x=\sum_\gs \go(x^\ttk_{\gs})\ot a_\gs,\ee with    
$ a_\gs\in S(\fh)$. Then define $\gk(x)= \sum_\gs  x^\ttk_{\gs}\ot\gk(a_\gs)\in S(\fb^-)$.  This carries more information than the usual symbol map for $U(\fb^-)$.
\\ \\
In one of our main results, Theorem \ref{zra} we obtain 
a closed formula
for the leading terms of the coefficients 
of  $\theta_{\gc,m}$ in terms of the leading terms of the coefficients of 
$\theta_{\gb,m}$, and three algebra maps.  We give a short summary.
\subsubsection{Algebras related to $S(\fh)$} \label{ta2} 
Fix an integer $m$. 
By \eqref{vat} the simple reflection $s_\ga$ maps $
{H}_{\gb, m}$ to ${H}_{\gc, m}$ %Then  $s_\ga   {H}_{\gb, m} =H_{\gc, m} $.  
Then  ${s_\ga} $ induces maps $P_{\ga} :S(\fh)\lra S(\fh)$ and 
$\cO(H_{\gc, m}) \lra \cO(H_{\gb, m}) $
such that $P_{\ga}  f(\mu) =f({s_\ga}  \cdot \mu)$.  
 Also 
  $\cO(\fh^*) = S(\fh)$ and we have a commutative diagram
\Bc
\xymatrix{&&&
 &S(\fh)\ar@{>}[d]\ar@{>}[rr]^{P_{\ga} }
&%S(\fh)\ar@{>}[dr]\ar@{>}[dl]
&S(\fh)\ar@{>}[d]\\%\ar@{-}[dr]\\
&&&&\cO(H_{\gc, m}) \ar@{>}[rr]
&&\;
\cO(H_{\gb, m})%\ar@{>}[ul]
&
}\Ec
where the downward arrows are restrictions of functions to a hyperplane. 
The map   $P_\ga$
\index{The map $P_\ga$}
is an  important ingredient in both Theorem \ref{cgw} and Theorem \ref{zra}.

\subsubsection{Algebras related to $S(\fn^-)$ and the algebras $R(\gt,\gs)$} \label{ta1} 
\noi Fix  $\gt, \gs\in \{\gb,\gc\}$.  The 
{\it $\gt$-Šapovalov algebra} Š$_\gt$ \index{$\gt$-Šapovalov algebra} is defined by taking the direct sum weight spaces over all nonnegative integers $m$, of weight $-m\gt$ in  the localization $S(\fn^-)_e$, 
see  Section \ref{ve} for more details. 
 The  
{\it small $\gt$-Šapovalov algebra}   $\mbox{Š}^0_{\gt}$
\index{$\gt$-Šapovalov algebra!small}  is the graded subalgebra of $\mbox{Š}_{\gt}$ generated by the degree 1 part $\mbox{Š}_{\gt}(1)$. 
%If $\gb$ is  isotropic then $\mbox{Š}^0_\gb$ and $\mbox{Š}_\gb$ are finite dimensional.
 The analogs 
$B_S, \nabla^t_S$ of $B_U$ and $\nabla^t_U$ are isomorphisms of
 graded  algebras
\be\label{781}\mbox{Š}_\gb  \stackrel{\nabla^t_S}{\lra} \mbox{Š}_\gb \stackrel{B}{\lra} \mbox{Š}_\gc.\ee
 Here $\nabla^t_S = \exp(td_ -)$ where $d_ -$ is a locally nilpotent derivation of
$S(\fn^-)_e$  that preserves  weight spaces.  
The map $B_S$ on  $\mbox{Š}_\gb (m)$  is multiplication by $e_{-\ga}^{mq}$.  
%The problem with $B_S$ is
 Note that the multiplication takes place in $S(\fn^-)_e$. A rigorous definition of $B_S$ is given in Equation \eqref{ma4}.  The maps in  \eqref{781} restrict to isomorphisms of small Šapovalov algebras. 
The $\gt$-Šapovalov algebra is the largest algebra such that \eqref{ma4} holds while $\mbox{Š}^0_{\gt}$ is the smallest.
\\ \\
Next set  
 $R(\gt,\gs) =\mbox{Š}_{\gt}^0\ot \cO(H_{\gs,0}).$
%\index{Graded ring $R(\gt,\gs)$}
 If $\gt$ is non-isotropic,  then by shifting the origin in $\fh^*$ we can make  
$R(\gt,\gs) $ into a graded algebra 
with  degree $m$ component  $\mbox{Š}^0_{\gt}(m)\ot \cO(H_{\gs,m}),$ see Lemma \ref{td5}. \index{Algebras $R(\gt,\gs)$}
\noi Combining the  above  maps we obtain algebra isomorphisms.
%\be \label{zra3} 1\ot P_\ga: R\gc,\gc\lra  R\gc,\gb\ee and 
\be B_S\ci\nabla_S^t\ot 1: R(\gb,\gb) \lra  
R(\gc,\gb).\nn\ee 

\be \label{zra4}  1\ot P_\ga: R(\gc,\gc) \lra  
R(\gc,\gb).\nn\ee 




\bt \label{783} If $t $ is as in  
\eqref{121d}, 
then the leading terms of the coefficients of the Šapovalov elements $\theta_{\gamma,m}=\sum_{\pi }  x_{\pi} c_\pi$ and $\theta_{\gb,m}
=\sum_{\gs }  x_{\pi} a_\gs$ are related by
\by 
\label{178} \sum_{\pi \in {\bf P}(m\gc)}  (1\ot P_\ga)( x_{\pi} \gk(c_\pi))
&=& \sum_{\gs \in {{\bf P}}(m\gb)} (B_S\ci\nabla_S^t\ot 1)(x_{\gs}  \gk(a_\gs)). 
\ey 
\et\noi
Theorem \ref{783} is a special case of  Theorem \ref{zra} which  applies to more general Lie superalgebras and uses $\Z$-forms of all the algebras involved.
\\ \\Rearranging terms in \eqref{178} we  have an isomorphism, see \eqref{2W4}  
\be \label{eeh} f: =
(1\ot P_\ga)\ci ((B_S\ci\nabla_S^t)\ot 1): R(\gb,\gb)\lra  R(\gc,\gc).\ee  such that 
$f(\gk({\theta}_{\gb,m}))=\gk({\theta}_{\gc,m}).$ 
\\ \\  
This paper is organized so that the results on coefficients can be presented as soon as possible.
\noi The next Section contains some preliminaries. After stating the key Lemma \ref{11768},  we sketch the induction step in the construction of \v Sapovalov elements. We also explain how to zoom in on the induction step, which allows us to greatly simplify notation at a cost of losing information from previous steps. 
The main results concerning coefficients are stated in Section \ref{PA.2} and proved in Section \ref{1s.5}. The proofs depend on a rather subtle cancellation property which is illustrated in Subsection \ref{Type C} in the cases of $\fsp(6)$ and $\osp(2,4)$. In 
Section \ref{g1}, apart from the results described above, we use convex geometry to show that 
the algebras Š$_\gc$  and Š$_\gb$  are finitely generated %normal and Cohen-Macaulay 
(provided $\fg $ is finite dimensional). We also study module categories for \v Sapovalov   algebras. 
In Section \ref{GF} we interpret our results geometrically. The final Section gives some examples.
\\ \\%The question of when there is an irreducinle map between Vermas will be discussed elsewhere.  
I would like to thank 
Kevin Coulembier,   Maria Gorelik, 
Volodymyr Mazorchuk, Dmitriy Rumynin,  Michel van den Bergh and the referee for some helpful correspondence.

\section{Preliminaries} \label{sss7.1}
%\noi It would be best to work over $\Z$, as in \cite{H} Chapters 25 and 26.\\ \\
%Before we state the main results, we introduce some notation.
 Throughout $[n]$ denotes the set of integers 
$\{1,2,\ldots,n\}$ and 
$\ttk$ is a field of characteristic zero. Let $A$ be a $\ttk$-algebra.  
A $\Z$-{\it  form} of $A$ 
 is a subring $A_\Z $  such that 
$A_\Z\ot_\Z \ttk =A$. 
If $A_\Z $  is $\Z$-{\  form} of $A$ 
 and $C$ is a commutative ring, we set $A_C =A_\Z\ot_\Z C$. 


\subsection{The Lie superalgebras} \label{gsi}For $\ttk$  algebraically closed, we make the assumptions below about the Lie superalgebra  $\fg$.
\bi
\itema  $\fg=\fg(A,\gt)$ %is a KM\ff{Kac-Moody Lie superalgebra 
\index{KM Lie superalgebra $\fg(A,\gt)$} is a KM\ff{Kac-Moody Lie superalgebra. These algebras are also known as contragredient Lie superalgebras, \cite{Kac1}, \cite{M101}  Chapter 5} Lie superalgebra with symmetrizable Cartan matrix $A$
%itemb No root string for a nonisotropic simple root has length greater than 2.
\itemb $\fg$ is integrable, that is
the adjoint action of any nonisotropic (positive or negative) root vector is locally nilpotent.
\itemc The enveloping algbra $U(\fg)$ \index{Integral forms!$U_\Z(\fg)$} has a $\Z$-form $U_\Z(\fg)$ that is stable under the divided power adjoint action of (positive and negative) root vectors for all simple nonisotropic roots.
\ei The  
Lie superalgebras satisfying (a)-(b) have been classified, \cite{S1}. They include all finite dimension simple Lie algebras and all basic classical simple Lie superalgebras.  However in place of $\fpsl(n|n)$ we use $\fgl(n|n)$.
\\ \\
If $\fg$
 is a simple Lie algebra, then $U_\Z(\fg)$ was constructed by Kostant based on work of Chevalley, see \cite{H}   Chapters 25-27. Fioresi and Gavarini \cite{FG} showed that hypothesis (c) also holds for classical Lie superalgebras of type A-D using an analogous construction. 
For KM Lie algebras given by a GCM, hypothesis (c) holds by \cite{Mq} Chapter 7. 
 For the Lie superalgebra $\fg=D(2,1;\ga)$ with $\ga$ irrational, there is obviously no $\Z$-form of $U(\fb^-)$.
\\ \\
Let 
 $\fg(A,\gt)$ be a KM Lie superalgebra with Cartan subalgebra $\fh$, and set of simple roots $\Pi$.
\v Sapovalov elements need not exist for Lie superalgebras see Example \ref{2.5}, so we make the following definition. Let $W_\Pi$ be the subgroups of the Weyl group $W$ generated by all reflections corresponding to roots 
nonisotropic roots in $\Pi$. We say $\gb$ is a {\it real root}\index{Real root} if it is in the $W_\Pi$ orbit of $\Pi$. This is different from the definition in \cite{S1}.\\ \\
If  (a)-(c) hold, we show that for real roots \v Sapovalov elements and their evaluations can also be constructed inside   the algebras $U_\Z(\fb^-)$ and $U_\Z(\fn^-)$.  In cases where (a), (b) hold but (c) does not, some results still work over $\ttk$. 
\\ \\ The assumption that $\ttk$ is algebraically closed is used only to limit the algebras considered. % For example we do not consider (nonsplit) real forms.


\subsection{Triangular decomposition and bilinear form}
 Here we work over $\ttk.$ 
 \index{Simple roots and positive roots $\Pi$, $\Delta^{+}$}
Let $\Delta^{+}$ be the set of positive roots  containing $\Pi$,
 and %-
\be \mathfrak{g} = \mathfrak{n}^- \oplus \mathfrak{h}
\oplus \mathfrak{n}^+\nn\ee \index{Triangular decomposition of $\mathfrak{g}$}
  the corresponding triangular decomposition  of $\fg$. We use the Borel subalgebras $\mathfrak{b} =  \mathfrak{h}
\oplus \mathfrak{n}^+$ and
$\fb^- =
\mathfrak{n}^- \oplus \mathfrak{h}$.
  Verma modules are  induced from $\mathfrak{b}$.  We do not use other Borel subalgebras. 
%Denote the unique simple factor of $M(\gl)$ by $L(\gl)$.
\\ \\
%\noi By \cite{M101} Theorem 5.4.1 and (5.4.3), 
Since $A$ is symmetrizable there is a nondegenerate invariant   bilinear form on $\fg$ \index{Non-degenerate bilinear form on $\fg$} which induces   an invariant  nondegenerate invariant symmetric bilinear form $(\;,\;)$ on $\fh^*$. For all $\ga \in \fh^*$, let $h_\ga \in \fh$ be the unique element such that \be \label{hdef}(\ga,\gb) = \gb(h_\ga) \mbox{ for all }\gb \in \fh^*.\ee Then for all $\alpha
\in \Delta^+$, choose elements $e_{\pm \alpha} \in
\mathfrak{g}^{\pm \alpha}$
 such that

\be \label{1eeh} 
[e_{\alpha}, e_{-\alpha}] = h_{\alpha}.\ee
It follows that if
$v_\mu$ is a highest weight vector of weight $\mu$ then
\be \label{ebl} e_\ga e_{-\ga} v_\mu = h_\ga v_\mu =(\mu, \ga)v_\mu . \ee
\noi 

\subsection{Weyl group}
Suppose $\Pi_{\rm nonisotropic}$  is the set of nonisotropic  simple roots.
\index{Set of nonisotropic  simple roots  $\Pi_{\rm nonisotropic}$}
  \noi
%Following \cite{S1} Defintion 3.5 and Section 10, we call an even root $\ga$ {\it principal} if there is  a Borel subalgebra $\fb'$ obtained from a $\fb$ by a sequence of odd reflections such that $\ga/2$ is a simple root of $\fb'$.
 For a nonisotropic root $\ga,$ we set $\alpha^\vee = 2\alpha /
(\alpha, \alpha)$, and denote the reflection corresponding to $\alpha$ by $
s_\ga$.      Also define 
\be \label{efh} t_{\alpha} = 2h_{\alpha}/(\alpha, \alpha).\ee For all $\gb\in \fh^*$, we have 
%\begin{equation} \label{27a}
$t_{\alpha} (\gb) = ({\alpha}^\vee,\gb).$
The  Weyl group  $W$ is the subgroup of $GL({\fh}^*)$
generated by all reflections.  For $u \in W$ set
 \be %\label{nu} 
N(u) = \{ \alpha \in \Delta_0^+ | u \alpha < 0 \},\qquad \ell(u) = |N(u)|.\nn\ee
We need the following well-known fact.
\bl
If $w  = s_\ga u$ with  $\ell(w)>\ell(u)$ and $\ga$ is a nonisotropic simple root, then we have a disjoint union
\begin{eqnarray} \label{Nw}
N(w^{-1}) = s_\ga N(u^{-1}) \ds \{\ga\}.
\end{eqnarray}
\el\noi 
%\end{document
Integrability has some important consequences. First set 
 \index{Special sets of roots ${\overline{\Delta}}^+_{0}, {\overline{\Delta}}^+_{1}$}
\be % \label{Kacroot}  
{\overline{\Delta}}^+_{0} = \{ \alpha \in \Delta^+_{0} | \alpha
/ 2 \not\in \Delta^+_{1}\}, \quad {\overline{\Delta}}^+_{1}= \{ \alpha
\in \Delta^+_{1}|2\alpha \not\in \Delta^+_{0}\}.\nn\ee
If $\ga$ is a nonisotropic root, then $\exp e_{\ga}\exp e_{-\ga}\exp e_{\ga}$ is a linear isomorphism  sending $\fg^\nu $ to  $\fg^{s_\ga\nu} $.
Thus $\dim \fg^\nu = \dim\fg^{s_\ga\nu} $.%We say $\gb$ is a real root if $w\gb\in \Pi$ for some $w\in W.$
\bl \label{xyz} If $\gb$ is a real root then
\bi \itema $\dim \fg^\gb =1$\ei
Suppose $k\gb$ is a root of $\fg$ for $k$ positive, then 
\bi  
\itemb If $\gb\in
{\overline{\Delta}}^+_{0}  \cup {\overline{\Delta}}^+_{1}$ then $k=1$.
\itemc If $\gb$ is odd non-isotropic, then  $k =1$ or $2$.
\ei
\el
\bpf Since this is true for $\gb$ simple, the result follows from the preceding remarks.\epf

\subsection{Symmetrization and derivations} \label{ss6.4} 
%We summarize some results from \cite{M101} Chapter 6.
Let 
 $\{ U_n \}$ be the standard filtration on $U = U(\fn^-)$ and  
     let $S^n(\fn^-)$ be the graded component of the symmetric algebra     $S(\fn^-)$  of degree $n$. 
There is a linear  map  $\omega_n: S^n(\fn^-) \longrightarrow U$ defined  by
\[ \omega_n(x_1x_2 \ldots x_n)   = (1/n!)\sum_{\sigma  \in \mathcal{S}_n}
\gamma(\pi(x),\sigma ) x_{\sigma(1)}x_{\sigma(2)}\ldots
x_{\sigma(n)} \] where $\gamma(\pi(x),\sigma) =\pm 1$ is as in \cite{M101} Equation $(A.2.36)$. If $\mathbb{S}^n(\fn^-)$ denotes the  image of 
 $S^n(\fn^-)$ in $U$, we have      $U_n = U_{n-1} \oplus \mathbb{S}^n(\fn^-)$.   Let  $$\go:S(\fn^-)=\bigoplus_{n\ge 0} S^n(\fn^-) \longrightarrow \bigoplus_{n\ge 0}  \mathbb{S}^n(\fn^-) =U(\fn^-)$$ 
\index{Symmetrization map $\go$}
be  the isomorphism of
vector spaces which coincides with
 $\omega_n$ on $S^n(\fn^-)$. 
\subsubsection{Derivations} \label{taa} 
 Let $A$ be an associative algebra For  $e, a \in A$
     and all $r \in \mathbb{N} $ we have,
\be \label{1cow} e^{r}a = \sum^r_{i=0} \left( \begin{array}{c}
                r \\
                i \end{array}\right) ((\ad e)^i a)e^{r - i} .\ee
\noi Here we interpret $(\ad e)a $ as $ea-ae$.  The following consequence is well-known.% \cite {M3}.
  
\bc \label{1oreset} If $e$ is a locally {\rm ad}-nilpotent nonzero divisor in $A$,  then
$\{e^n|n \in \mathbb{N}\} $
is an Ore set in $A$. We write $A_e$ for the resulting Ore localization.
\ec
\noi \noi If $\ga \in \Pi_{\rm nonisotropic}$, and  $e = e_{-\ga}.$ Then $e$ satisfies the hypothesis in 
$U = U(\fn^-)$.
\\ \\If $\ga $ is a nonisotropic root, let
$\fl$ \index{Lie superalgebra $\fl$} be the subalgebra of $\fg$ generated by $e_{\pm \ga}$ and  $h_\ga =[ e_{\ga}, e_{-\ga}]$. Thus $\fl \cong \fsl(2)$ for  $\ga$ even or $\fl \cong \osp(1|2)$ if $\ga$ is odd. 
 Let $\fr$ be the ideal of $\fn^-$\index{An ideal of $\fr$ of $\fn^-$}
formed by negative weight vectors with weight different from $-\ga$. Then $\fn^-= \fr\op \ttk e_{-\ga}$.
\\ \\
There is a simple trick that allows us to treat the cases where $\ga$ is an even root and an odd nonisotropic root uniformly.
If $e$ is homogeneous in $\fg$ 
then $\pd=\ad  e$ is the  superderivation of $\fg$ defined by $(\ad  e)(x) = [e,x]$.  
Since  $U$ is  a $\Z_2$ graded algebra it has  an automorphism $s$  given by 
$$s(u_0 +u_1) = u_0 - u_1 \mbox{ for } u_i \in U_i.$$
If  $e$ is  odd, then  $\ado e=-s\ci\ad e\ci s$  is a derivation of $U$. 
Note that  $\ado e$ and $\ad e$ have the same restriction to $\fg$, compare the twisted adjoint action of Gorelik, \cite{G3}. \index{Twisted adjoint action $\ado e$}%It follows that $D$ preserves the standard filtration on $A$.
%\subsection{Adjoint action of $\fl$} \label{PA.1} 
\\ \\
Denote the extensions of $\ad e_{\pm \ga}$ or
$\ado e_{\pm \ga}$  (for $\ga$ even or odd respectively) to derivations 
$\pd_S^{\pm}
,  \pd_U^{\pm}$ of $S(\fr^-)$, $ U(\fr^-)$. %Previously $\pd_S^{-},  \pd_U^{-}$ were written $\pd_S,  \pd_U$. 
Then extend these maps to derivations of $S(\fn^-)$  or $U(\fn^-)$ that vanish on $e_{-\ga}$. Observe that if $x\in U(\fn^-)$, then $ \pd_U x 
=ex-xe$ for any parities. We often write $\pd_S,  \pd_U$ in place of   $\pd_S^{-},  \pd_U^{-}$ since they are used far more frequently than the other two derivations.

 \bl \label{npr}
The map $\omega:S(\fn^-)
     \longrightarrow U(\fn^-)$ is $\pd$  equivariant, that is  $\go\pd_S(u) = \pd_U\go(u)$.
\el \bpf This is well-known, \cite{M101} Theorem 6.4.4.
%Then apply the Theorem with $\fl$ in place of $\fn^-.$  
\epf\noi   
\subsection{Partitions}%$\upsilon$
\noi 
Set $Q^+=\sum_{\gz\in \Pi} \N\gz$.
If $\eta \in Q^+$, a {\it
partition} \index{Partitions} of $\eta$ is a map
$\pi: \Delta^+  \longrightarrow
\mathbb{N} $ such that
$\pi(\upsilon) = 0$ or $1$ for all isotropic roots $\upsilon$,
 $\pi(\upsilon) = 0$ for all even roots $\upsilon$ such that $\upsilon/2$ is a root, and
\be  %\label{suma}
\sum_{\upsilon \in {{ \Delta^+}}} \pi(\upsilon)\upsilon = \eta\nn.\ee
For $\eta \in Q^+$, we denote by $\bf{{P}(\eta)}$ the set of partitions of $\eta$. If $\ga$ is a positive isotropic root, set ${\bf{P}}_\ga(\eta)=\{ \pi\in {\bf{P}}(\eta)|\pi(\ga)=0\}$, $\bf{{p}(\eta)}=|\bf{{P}(\eta)}|$ and   $\bf{{p}_\ga(\eta)}=|\bf{{P}_\ga(\eta)}|.$ \index{Partitions!sets of ${{\bf P}}, \bf{P}_\ga$} \index{Partitions!number of $\bf{{p}(\eta)},\bf{{p}_\ga(\eta)}$}
%If $\gs, \pi$ are partitions we say that $\gs + \pi$ is a partition if$\gs(\upsilon)+\pi(\upsilon) \le 1$ for all isotropic roots  $\upsilon$. In this case   $\gs + \pi$ is defined by$(\gs + \pi)(\upsilon) = \gs(\upsilon)+\pi(\upsilon) $  for all positive roots  $\upsilon$.
  Fix a Chevalley basis $\{e_{-\upsilon}| \upsilon\in \Gd^+\}$, $\{h_\upsilon| \upsilon \in \Pi\}$ for $\fn^-, \fh$ respectively. 
For each positive root $\upsilon$ we have a root vector 
$e_{-\upsilon}$. Let   $x_{\upsilon}$ be another copy of 
$e_{-\upsilon}$ but inside $S_\Z(\fn^-)$. 
Fix an order on the set $ \Delta^+$, and for $\pi$ a partition,
set 
\be \label{negpar} e_{-\pi} = \prod_{\upsilon \in \Delta^+} \frac{e^{\pi (\upsilon)}_{-\upsilon}}{{\pi (\upsilon)!}}, \quad x_{\pi} = \prod_{\upsilon \in \Delta^+} \frac{x^{\pi (\upsilon)}_{\upsilon}}{{\pi (\upsilon)!}}\ee \index{Partitions!elements $e_{-\pi},  x_{\pi}$
associated to}\noi 
the  products being taken with respect to this order.  Fix $\ga\in \Pi_{\rm nonisotropic}$.  We say the order on $ \Delta^+$ is an {\it $\ga$-order} \index{$\ga$-order} if the terms involving 
$e_{-\ga}$ and 
$x_{\ga}$
occur last in 
\eqref{negpar}. The elements $e_{- \pi}$ and  $x_{\pi} $ with $\pi \in \bf{{P}}(\eta)$ 
form a $\Z$-basis of
$U_\Z(\mathfrak{n}^-)^{- \eta}$ and $S_\Z(\mathfrak{n}^-)^{- \eta}$
\index{Integral forms!$U_\Z(\mathfrak{n}^-), S_\Z(\mathfrak{n}^-)$}
 respectively.   
To work over $\ttk$ we add a superscript and define 
\be \label{kpar} 
e^\ttk_{-\pi} = \prod_{\upsilon \in \Delta^+} {e^{\pi (\upsilon)}_{-\upsilon}}, \quad x^\ttk_{\pi} = \prod_{\upsilon \in \Delta^+} x^{\pi (\upsilon)}_{\upsilon}.\ee 
If $\pi \in \bf{{P}(\eta)}$  define  $|\pi| = \sum_{\upsilon \in \Delta^+} \pi(\upsilon).$ Then 
 $|\pi|$ is the degree of $e_{-\pi}$ in the standard filtration.  It is useful to extend this notation to arbitrary elements of 
$U(\mathfrak{n}^-)$. Thus if $x\in 
U_n(\mathfrak{n}^-) \backslash 
U_{n-1}(\mathfrak{n}^-)$ we set $|x|=n.$
\\ \\
There are two special partitions of $m\gc$. Let
$\pi^0 \in {{\bf P}}(m\gc)$ be the unique partition of $m\gc$ such
that $\pi^0(\upsilon) = 0$ if $\upsilon \in \Delta^+ \backslash \Pi .$ The partition $m\pi^{\gc}$ of $m\gc$ is given by $m\pi^{\gc}(\gc)=m,$ and $m\pi^\gc(\upsilon)= 0$ for all positive roots $\upsilon$ different from $\gc.$
The partition $m\gs^{\gb}$ of $m\gb$ is defined similarly. 
This means that $e_{-m\pi^{\gc}} =e_{-\gc}^m/m!$ 
and  $e_{-m\gs^{\gb}} =e_{-\gb}^m/m!$. 
\\ \\
If { $\pi \in {{\bf P}}(m\gc), \pi \neq
m\pi^{\gc}$ and  $|\pi|\le m$} we call $\pi$   a {\it rogue partition of }$
{m\pi^{\gc}}.$ \index{Partitions!rogue}

\subsection{Coefficients of \v Sapovalov  elements} \label{s.2}
Fix $\gr\in \fh^*$ such that if $\ga$ is a simple root, then $(\gr,\ga) =(\ga,\ga)/2$. \index{The element $\gr$} The dot action of $W$ on $\fh^*$ is defined by $w\cdot \gl = w(\gl+\gr) -\gr$. \index{The element $\gr$!  dot action}
Fix a positive root $\gc$,  a  positive   integer $m$ and  an  
order on $ \Delta^+$. If $\gc$ is isotropic, assume $m=1$, and if $\gc$ is odd nonisotropic, assume that $m$ is odd.
We say that $\gth = \theta_{\gamma,m}\in U({\mathfrak n}^{-} )^{- m\gc} \ot S(\fh)$ is a \index{\v Sapovalov element $\theta_{\gamma,m}$}
{\it \v Sapovalov element for the pair} $(\gc,m)$ if it has the form
\be \label{rat}
\theta = \sum_{\pi \in {{\bf P}}(m\gc)} e^\ttk_{-\pi} 
 c_{\pi},\ee where 
\be \label{rbt} c_{\pi} \in S({\mathfrak h}), \; c_{\pi^0} = 1,\nn\ee  and
\be %\label{boo} 
e_{\upsilon} \theta \in U({\mathfrak g})(h_{\gc} +{\gr}(h_\gc)
-m(\gc,\gc)/2)+U({\mathfrak g}){\mathfrak n}^+ , \; \rm{ for \; all }\;\upsilon \in \Pi. \nn\ee
We call the $c_{\pi}$ in \eqref{rat} the {\it coefficients} \index{\v Sapovalov element $\theta_{\gamma,m}$!coefficients of} of $\gth$.  
 For a simple Lie algebra, the existence of such elements was shown by \v Sapovalov, \cite{Sh} Lemma 1.
%The normalization condition $c_{\pi^0}=1$ ensures that $\theta_{\gc, m}v_\gl$ is never zero.
%If $m=1$, to simplify notation  set $\mathcal{H}_{\gamma} = \mathcal{H}_{\gamma,1},$ anddenote a  \v Sapovalov element for the pair $(\gc, 1)$ by $\gth_{\gamma}$.  
\\ \\
We evaluate $\theta_{\gamma,m}$ at points $\gl\in H_{\gc, m}$. %Writing   $\gk( c_{\pi}$ for the restriction of $c_{\pi}$ to $H_{\gamma,m}$ it is tempting to write \eqref{rat} in the form
Abusing notation we also write \eqref{rat} in the form $$
\theta = \sum_{\pi \in {{\bf P}}(m\gc)} e^\ttk_{-\pi} 
 c_{\pi}\in U({\mathfrak n}^{-} )^{- m\gc} \ot 
\cO( H_{\gc, m}).$$ Then the functions $c_{\pi}$ are uniquely determined.  
\bexa \label{2.5}{\rm \v Sapovalov elements need not exist for Lie superalgebras.  Let $\fg=\osp(3,2)$ This algebra has two systems of positive roots up conjugacy.  Namely
$$ \Pi_1 = \{\gep-\gd, \gd \}, \quad \Pi_2 = \{\gd-\gep, \gep \}.$$
Here  $\gep$ is an even root $\gd$ is odd nonisotropic and $\pm (\gep-\gd)$ are isotropic. Suppose we try to find a \v Sapovalov element for the pair $(\gep,m)$ using $\Pi_1$.  The root vectors $
e_{\gep-\gd}$ and $e_{-\gep}$ are given by the KM construction and we define $e_{-\gd}= [e_{\gep-\gd}, e_{-\gep}]$.  We choose $e_{\gd-\gep}$ so that \eqref{1eeh} holds. 
Let $M(\gl)$ be a Verma module constructed using this system. Since $\gep$ is not a simple root of $\Pi_1$,  we use an odd reflection to pass to  $\Pi_2$, multiply by  $ e_{-\gep}^m$ and use another odd reflection to return to  $\Pi_1$. This yields  
\by e_{\gep-\gd} e_{-\gep}^m e_{\gd-\gep}v_\gl &=& ( e_{-\gep}^m + me_{-\gd} e_{-\gep}^{m-1})e_{\gd-\gep}v_\gl\nn\\
&=& ((\gl,{\gep-\gd}) e_{-\gep}^m  + me_{-\gd} e_{-\gep}^{m-1}e_{\gd-\gep})v_\gl\nn \ey This is a highest weight vector, but since the coefficient of $ e_{-\gep}^m$ is not constant, there is no  \v Sapovalov element.
}\eexa
 % and let $W_\Pi_{\rm nonisotropic}$ (resp. $W_\Pi_{\rm even}$) bethe subgroup of $W_\Pi$ generated by the reflections $s_\ga,$ where $\ga\in \Pi_{\rm nonisotropic}$ (resp. $\ga \in \Pi_{\rm even}$)
\subsection{Evaluation and a key  Lemma} \label{eva} Let $R$ be a $\ttk$-algebra. 
If  $\gl \in {\mathfrak h}^* $ we define  {\it evaluation at}  \index{Evaluation} $\gl$ to be the
map 
\be %\label{eval}
\varepsilon^\gl:R\otimes S({\mathfrak h})  \longrightarrow R,\;  \quad
\sum_{i} a_i \otimes b_i \longrightarrow \sum_{i} a_i  b_i(\gl).\nn\ee 
%For now we only use this notation with $R= U({\mathfrak n}^-)$ or $U({\mathfrak n}^-)_e$ and we identify  $U({\mathfrak n}^-)\otimes S({\mathfrak h})$ with $U({\mathfrak b}^-).$
\noi 
\noi We construct \v Sapovalov elements  inductively using the next Lemma. Fix an $\ga$-order on $ \Delta^+$. \index{Key lemma of \v Sapovalov} 
\begin{lemma}\label{11768}
Suppose $\gb$ is a real root, $\mu \in H_{\gb,m}, \;\alpha \in\Pi_{\rm nonisotropic} \cap A(\mu)$ and set
\be \label{121d}\gl = s_\ga\cdot\mu,\;\;\gc = s_\alpha\gb,\;\;p = (\mu, \alpha^\vee),\;\;q = (\gc, \alpha^\vee)=-(\gb, \alpha^\vee).\ee
Assume that $ q, p  \in \mathbb{N}\backslash \{0\}$, $t = mq+p$ and \begin{itemize}
\itema  $\theta' \in
U({\mathfrak n}^-)^{-m\gb}$ is such that $v = \theta'v_\mu \in
M(\mu)$ is a highest weight vector.
\itemb  If  $\ga$ is odd, 
then $q = 2$ and $p$ is odd.
%\itemc If $\alpha$ satisfies \eqref{1bog} then $q=1.$
\end{itemize}
\noi Then there is a unique $\theta \in U({\mathfrak n}^-)^{-m\gc}$ such that
for all $\mu \in H_{\gb,m}$
\begin{equation} \label{121nd}
e^{p + mq}_{- \alpha}\theta' = \theta e^p_{- \alpha}.
\end{equation}  
\end{lemma} 
\bpf If $\ga$ is even this  is well-known, see for example  \cite{H2} Section 4.13 or \cite{M101} Lemma 9.4.3. \ff{For the corrected proof see: https://www.ams.org/publications/authors/books/postpub/gsm-131-Errata.pdf} 
The proof  uses the representation theory of $\fsl(2)$.
If $\ga$ is odd, non-isotropic we use the representation theory of $\osp(1,2)$. We can choose elements  $x = e_{\ga},\; y= e_{-\ga}$, and  $h = [x, y]= h_{\ga}$ in $\fg$ such that  $[h,x] =x$, $[h,y] =-y$ and  $\fk=\span \{x, y, h \}\cong \osp(1,2)$.   
By \cite{M101} Lemma 8.3.2,  $[x, y]= (x,y)h_{\ga}$, so $(x,y)=1$.  Therefore
$$1= ([h,x],y) = ([h,[x,y]) = (h,h) = (\ga,\ga)$$ so $\ga^\vee =2\ga$.   
 Under the given hypothesis we may identify $M(\gl)$ with the submodule of $M(\mu)$ generated by $v_\gl=y^pv_\mu$. 
The element $v$ from (a) has weight $\mu- m\gb$. 
Since $\ad y$ is locally nilpotent there is a positive integer $\ell$ such that 
$y^\ell v \in M(\gl),$ and we may assume $\ell > p +2m$ and  $\ell$ is odd.  %Let $\gk' = m\gb$. 
Then $(\mu- m\gb,\ga) = (p-1)/2 +m:=a$, so $hv = av.$ 
We compute that
\[ x^2y^{\ell}v =  \frac{\ell -1}{4}(\ell -p-2m)y^{\ell-2}v \in M(\gl).
\] Since  $\ell > p +2m$, we deduce that $y^{\ell-2}v \in M(\gl).$ 
Repeating this argument leads to $y^{p+2m}\theta'v_\mu \in M(\gl)= U({\mathfrak n}^-)y^pv_\mu.$ This shows the existence of $\gth$. Uniqueness follows since $y^p$ is not a zero divisor in  $U({\mathfrak n}^-)$.\epf
\noi %Equation \eqref{121nd} is the basis for the proof of many properties of \v Sapovalov elements
Set $e = e_{- \alpha}$.
If we are given $\theta_{\gb,m}$ and apply the Lemma with 
$\gth' = \theta_{\gb,m}(\mu)$  for $\mu$ in a dense subset of $H_{\gb,m}$ we obtain the \v Sapovalov element
$\theta_{\gamma,m}$ and  the relation

\be \label{12nd} 
e^t\theta_{\gb,m}(\mu) =\theta_{\gamma, m}(s_\ga\cdot\mu)e^p  % 
\ee
for all $\mu\in H_{\gb,m}.$ For $t\in\Z$, define $\Phi_t$ 
\index{Conjugation operator $\Phi_t$} by 
$\Phi_t(x) = e^txe^{-t}$.  If $t = mq+p$, as in the Lemma, we can rewrite \eqref{12nd} in the form
\be \label{12md} 
\Phi_t(\theta_{\gb,m}(\mu)  )e^{mq} =\theta_{\gamma, m}(s_\ga\cdot\mu)  % \nabla^{t}_U( \theta_{\gb,m})
.\ee  Note that $\gb$ is isotropic iff $\gc$ is. {\it We assume from now on that $\gb$ and  $\gc$ are nonisotropic unless otherwise stated.} A real root can be $\gb$ isotropic or nonisotropic.
\subsection{Deformed multiplication and zooming in}  \label{ma3}  
\subsubsection{Deformed multiplication}   
\noi  Set  $\gc =s_\ga \gb$ and $\varpi = \frac{\ga}{(\ga,\gb^\vee)} = -\frac{\ga}{(\ga,\gc^\vee)}.$ 
%We use multiplication in $S(\fh)$ and move the origin in $\fh^*$ to define certain algebra maps. % in the Lemma below.
 Fix  $\gs\in \{\gb,\gc\}$. 
%Note that $(\mu + m{\varpi},\gb^\vee) =m$ if $\mu\in H_{\gb, 0}$.
The map \be \label{jjj} y_m^*:H_{\gs, 0} \lra H_{\gs, m}.\ee
 given by 
$\mu\lra \mu \pm m\varpi$ 
where the sign is $+$ (resp. $-$) if $\gs = \gb$ (resp. $\gs = \gc$)
induces a map $y_{m}:\cO(H_{\gs,m})\lra \cO(H_{\gs,0})$.
%Below, we assume $\gs = \gb$.  If $\gs = \gc$ we should change the sign of all coefficients of $\varpi$. 
\bl \label{td5} There is an associative multiplication 
\be %\label{qi} 
\cO(H_{\gs,m}) \ot \cO(H_{\gs,n})\lra \cO(H_{\gs,m+n}).\nn\ee
where for $\gth\in \cO(H_{\gs,m}), \phi \in \cO(H_{\gs,n})$ we have 
\be %\label{q11}
(\gth\cdot \phi )(y_{m+n}\mu) =\gth(y_{m} \mu)\phi (y_{n}\mu),\nn\ee 
\el \noi
\subsubsection{Zooming in}  
\noi \index{Zooming in}
Suppose that  $\mu \in H_{\gb,m}$. % and gl=s_\ga\cdot\mu$. 
Let $I$ be the ideal of $S(\fh) $ consisting of functions that vanish on $ H_{\gb,m}$. Consider the line %$L =\ttk (\quad , \alpha^\vee)\subset \fh^*$
$L =\ttk (\quad , \alpha)\subset \fh^*$. Since $ H_{\gb,m}$ meets $L$ at the  single point $m\varpi$,  it follows that $S(\fh) = I\op \cO(L)$. 
Thus we have a map 
$k:S(\fh)  \lra \cO(L)$ which factors through $\cO( H_{\gb,0})$.  
With  $h_\alpha$ as defined in \eqref{hdef}, define $$T=k(2h_\ga/(\ga,\ga))+t+p.$$  Setting $T=t$,  the combined map sends $2h_\ga/(\ga,\ga)$ to $-p = (\gl,\ga^\vee) = \varepsilon^{\gl}(2h_\ga/(\ga,\ga))$ and the diagram below commutes.

\xymatrix{
&&&&\
\cO( H_{\gb,0})
 \ar@{->}[dr]_{\varepsilon^{\gl}}  \ar@{->}[rr]^{k}  && \ttk[T]
\ar@{->}^{T\lra t%(\gl, \alpha^\vee)
}[dl] &\\
&&&&&\ttk
} \noi
Taking  coefficients from $\ttk[T]$ lets us to zoom in on the induction step allows for a more conceptual approach.  However if we do this, we lose all information about previous steps in the induction. The map $f$ from \eqref{eeh} induces an isomorphism of graded $\ttk[T]$-algebras, with $\deg T =0$
$$\mbox{Š}_\gb^0[T] \lra \mbox{Š}_\gc^0[T].$$


\section{Main results on coefficients} \label{PA.2} 
%\section{Some results from \cite{M3}} \label{pqr} 
 We keep the notation of Lemma \ref{11768}, but use an arbitrary order on $ \Delta^+$. 
We need a Zariski dense  subset $\gL$ of  ${H}_{{\gb_0},m}$ such that  any $\gl\in \Lambda$ satisfies some integrality and positivity conditions, which allow us to apply Lemma \ref{11768} repeatedly. %To define $\gL$, 
For  $\gl\in\fh^*$, set
\[A(\lambda)_{0}  =  \{ \alpha \in \overline{\Delta}^+_{0} | (\lambda +\gr,
\alpha^\vee) \in \mathbb{N} \backslash \{0\} \}, \]
\[A(\lambda)_{1} = \{ \alpha \in \Delta^+_{1} \backslash \overline{\Delta}^+_{1} |(\lambda+\gr, \alpha^\vee ) \in 2\mathbb{N} + 1 \}, \]
%\end{eqnarray*}
and  \be \label{1tar}A(\lambda) = A(\lambda)_{0} \cup  A(\lambda)_{1},\ee then 
 $\gL$ is defined by %(see \eqref{rue} for the notation)
\index{Zariski dense  set $\gL$ }
\be %\label{1tar}
\Lambda =
\{\nu \in {H}_{{\gb_0},m}| \Pi_{\rm nonisotropic}\subseteq A(\nu)
\}
.\nn\ee  We also need  \be \label{Bdef} B(\lambda) = \{ \alpha \in \overline{\Delta}^+_{1} | (\lambda +\gr,
\alpha) = 0 \} . \ee \index{Sets of roots $A(\lambda) ,B(\lambda)$ }
We assume that \be %\label{iii} 
\gc = w{\gb_0} \mbox{ for a simple root } {\gb_0} \mbox{ and } w \in W_\Pi\nn\ee
%\be \label{ii} \gc = w{\gb_0} \mbox{ for a simple root } {\gb_0} \mbox{ and } w \in W_{\rm nonisotropic}.\ee 
and that $\ell(w)$ is minimal. For $\alpha \in N(w^{-1}),$ we define 
\be \label{iv} q(w,\ga) = (w\gb_0, \ga^\vee).\ee 
 If $\Pi $ is the set of simple roots, and $\gc = \sum_{\upsilon \in \Pi}
 a_\upsilon\upsilon,$ then the {\it height} \index{Height of a root} $\hgt \gc$ of $\gc$ is defined to be
$\hgt \gc =\sum_{\upsilon \in \Pi}
 a_\upsilon$.
Let ${\mathcal I}(H_{\gc, m})$ be the ideal of $S(\fh)$ consisting of  functions vanishing on ${H}_{\gamma,m}$, and $\cO (H_{\gc, m})= S(\fh)/{\mathcal I}(H_{\gc, m}).$
\noi
In the Theorem below we assume $\fg$ is a a Lie superalgebra 
satisfying (a)-(c) from Section \ref{gsi}.
We require $\gc$ to be a real  root such that 
\bi 
\itema
 If $\gc$ is isotropic, then $m=1.$
\itemb If $\gc$ is odd nonisotropic, then $m$ is odd and $q=2$. 
\ei  

\bt \label{1sag} For all
$\gl \in w\cdot  \Lambda$ we have
\begin{equation} %\label{12}
\theta_{\gamma,m}(\gl) = \sum_{\pi \in {{\bf P}}(m\gamma)} e^\ttk_{-\pi}c_{ \pi}(\lambda) \in
U_\Z({\mathfrak n}^-)^{- m\gamma} ,\nn
\end{equation}
where 
the coefficients $c_{\pi}(\gl) \in \ttk $ depend polynomially on $\gl \in
w\cdot  \Lambda$, $c_{\pi^0} = 1$ 
\be \label{10a}
 \deg  \; c_{\pi}(\gl) \leq m\hgt \gc - |\pi|,\nn\ee
and \ $c_{m\pi^{\gamma}}(\gl)$ is a polynomial function of
$\gl$ of degree $m(\hgt \gc - 1)$ with highest term equal to $c\prod_{\gs
\in N(w^{-1})}(\gl,\gs)^{mq(w,\gs)}$ for a nonzero constant $c$.
\et %IAN Eliminate $c$.
\noi 
If $\fg$ is a simple Lie algebra this result is due to Kumar and Letzter, see \cite{KL}. \\ \\
By induction we use an $\ga$-order to  write  \begin{equation} \label{1x7} \theta_{\gb,m} =
\sum_{\gs \in {{\bf P}}(m\gb)} e_{-\gs}a_{\gs}.
\end{equation} 
where the coefficients $a_{\gs}$ from \eqref{1x7} satisfy analogs of the conditions  in Theorem \ref{1sag}.
\\ \\ 
The proof of Theorem \ref{1sag} gives  the following closed fomula for all coefficients.
\bt \label{cot} With the same hypotheses as Theorem \ref{1sag}, the coefficients $c_{ \pi}(s_\ga\cdot\mu)$  of the Šapovalov element $\theta_{\gamma,m}$ as in \eqref{rat} are given by

\by \label{ceq}  c_{ \pi}(s_\ga\cdot\mu) %e^{t} \theta_{{\gb_0},m}(\mu)e^{N-t} =
%\sum_{\;j \geq 0}  \sum_{\gs \in {\overline{\bf P}}(m\gb+N\ga)}
=\sum_{j,\gs} 
 \q^{(j)} _{\gs}( e_{- \pi})
a_{\gs}(\mu).
\ey
\et \noi
\noi In \eqref{ceq} the $\q^{(j)} _{\gs}$ are some $\Z$-valued functions, which are defined using the adjoint action of an $\fsl(2)$-subalgebra of $\fg$ if $\ga$ is even, or an $\osp(1|2)$ subalgebra if $\ga$ is odd.
\\ \\
The integral form of 
$S(\mathfrak{h}) $ \index{Integral forms!$S_\Z(\mathfrak{h})$}
has $\Z$-basis consisting of all products $\prod_{\gz\in \Pi} \left(\begin{array}{c} h_\gz \\ c_\gz \end{array}\right)$, where $C= (\ldots c_\gz \ldots) \in \N^{\Pi}.$ Set 
$U_\Z(\fb^-)= U_\Z(\fn^-)\ot_\Z S_\Z(\mathfrak{h}).$
Note that no integrality properties are involved in the definition of a  Šapovalov element.  However   by \eqref{ceq} and induction 
\bc We have  $\theta_{\gc,m} \in U_\Z(\fb^-)$.
\ec
\bc \label{lter} In Theorem \ref{1sag}  $c_{m\pi^{\gamma}}$ is the unique coefficient of highest degree in $\theta_{\gc,m}.$
\ec
\bpf This follows easily from the given degree estimates, and the statements about the leading terms.\epf
%\subsection{The case where $\gc$ is isotropic.}
%\end{document}

\section{Proofs of Theorems   \ref{1sag} and \ref{cot}}\label{1s.5}

\subsection{The cancellation step and proof of Theorem \ref{cot} } \label{cs} Until Subsection \ref{zzw} we work with an $\ga$-order on $ \Delta^+$.  
\noi The main result of this Subsection, Theorem \ref{cot} relates
coefficients  of the Šapovalov elements  $\theta_{\gb,m}$ and $\theta_{\gamma,m}$. 
To prove the Theorem, we need some preparation.  
The functions $ \q^{(j)} _{\gs}$ that appear in  the statement  of Theorem \ref{cot} play an important role in the rest of the paper, see Equation \eqref{A89}. % and Lemma \ref{A8}. 
However before they can be defined we need a cancellation step. 
%In the next result we give a $\Z$-basis for the weight spaces of $U_\Z(\fn^-)_e$. Restricted to $U(\fn^-)$ it is different from the $\Z$-basis defined using \eqref{negpar}. 
%For $\gz \in \bf{{P}}(\eta)$ with $\gz(\ga) =k$, there is a unique $\pi \in{\bf{\overline{P}}}(\eta-k\ga). $such that $\pi(\upsilon) =\gz(\upsilon)$ for all $\upsilon\neq \ga$and we set \be \label{HYB}  e^\tth_{-\gz}=e_{-\pi} e^k \nn\ee
%Let ${\widehat{\bf P}}(\eta)$ be the set of pairs $(k,\pi)$ such that $k \in \mathbb{Z}, \pi \in {\overline{\bf P}}(\eta -k\ga)$ and $\pi(\ga) = 0.$  Then we have 
Let  $${\overline{\bf P}}(\eta)=\{\pi \in \bf{{P}}(\eta)| \pi(\ga)=0\}.$$
The elements $e_{- \pi}$ and  $x_{\pi} $ with  $\pi \in \bf{\overline{P}}(\eta)$  form a $\Z$-basis of 
$U_\Z(\mathfrak{r}^-)^{-\eta}$ and $S_\Z(\mathfrak{r}^-)^{-\eta}$.
%\index{Integral forms!$U_\Z(\mathfrak{r}^-), S_\Z(\mathfrak{r}^-)$}
 \index{Partitions!sets of ${\overline{\bf P}}$}

\bl \label{1uebasis}  \begin{itemize} \item[{}]
\itema The set $\{e_{-\pi} e^k|(k,\pi ) \in {\Z\ti\bf{\overline{P}}}(\eta-k\ga)\}$ forms a $\Z $-basis for the weight space $U_\Z(\fn^-)_e^{-\eta}.$
 \itemb If $u = \sum_{(k,\pi ) \in{\Z\ti\bf{\overline{P}}}(\eta-k\ga)} c_{(k,\pi )} e_{-\pi} e^k \in U_\Z(\fn^-)_e^{-\eta}$ with $c_{(k,\pi )} \in \ttk ,$ then $u \in U_\Z(\fn^-)_e
$ if and only if $c_{(k,\pi )} \neq 0$ implies $k \ge 0$.
\end{itemize}
 \noi  
\el \bpf    
Since    
 $U_\Z(\fn^-)_e =U_\Z(\fr^-)[e^{\pm 1}]$ is a twisted Laurent polynomial ring  the result follows.
 \epf
\noi With $\gz = (k,\pi )$ 
%$e_{-\pi} e^k$ 
as in (a), we set
$e^\tth_{-\gz}=e_{-\pi} e^k$.  This has weight $-\eta$.
We call the new basis the {\it hybrid} $\Z$-basis 
for $U(\fn^-)_e$, 
\index{Hybrid $\Z$-basis for  $U(\fn^-)_e$} 
since it involves divided powers of all negative root vectors other than $e$.  
This extends the usual notion of a partition: 
if  ${\gs \in {{\bf P}}(\eta)}$ 
and 
$\gs'\in {\bf{\overline{P}}}(\eta-\gs(\ga)\ga)$
satisfies  $\gs'(\upsilon) = \gs(\upsilon)$ for all $\upsilon\neq \ga,$ set 
$\gz =(\gs(\ga), \gs')$ and $e_{-\gs}^\tth = e_{-\gz}^\tth$. 
Then
 \be \label{JI} e_{-\gs}^\tth =\gs(\ga)!e_{-\gs}. \ee
% The next Lemma follows from the definition \eqref{negpar} of elements associated to a partition.
%\bl \label{JL}  If $\gz\in \overline{\bf P}(m\gb+N\ga)$ and $\gz(\ga)%\ge N-mq,$ there is a unique partition $\gz'\in \overline{\bf P}(m\gc)$ such that  $e_{- \gz'}e^{N-mq} = e_{ -\gz}$.  \el
We also set
\be \label{10b}e^\ttk_{-\gz}=e^\ttk_{-\pi} e^k.\ee
\bl \label{JL}  Suppose $\gz = (k,\pi )$ is as in (a) and $k\ge N-mq.$ Then  $\gz'  = (k+mq-N,\pi )$ satisfies 
 $e^\tth_{- \gz'}e^{N-mq} = e^\tth_{ -\gz}$ and $e^\tth_{- \gz'} \in U(\fn^-)$.  
\el 
\noi 
 

%Let ${\widehat{\bf P}}(\eta)$ be the set of pairs $(k,\pi)$ such that $k \in \mathbb{Z}, \pi \in {\overline{\bf P}}(\eta -k\ga)$ and $\pi(\ga) = 0.$  Then we have 
\noi
\subsubsection{The cancellation step} \label{pre} 
\bl \label{pre1} $\theta_{\gamma, m}(s_\ga\cdot\mu)\in U_\Z(\fn^-)^{-m\gc}$ for all   
$\mu\in H_{\gb,m}.$
\el
\bpf We use representation theory, so assume $t=p+mq$  as in Lemma \ref{11768}.
Assume  as in  Equation 
\eqref{1x7} that   $$\theta_{\gb,m} =
\sum_{\gs \in {{\bf P}}(m\gb)} e_{-\gs}a_{\gs} \in U_\Z(\fb^-).
$$
Fix $N$ large enough so that
 $\pd_U^{N+1} e_{-\gs} = 0,$ for all $\gs \in {{\bf P}}(m\gb).$ 
By Equations 
\eqref{1x7}, \eqref{JI} and \eqref{1cow}
\by \label{08x} e^{t} \theta_{\gb,m}(\mu) 
e^{N-t}
 & =& 
\sum_{\gs \in {{\bf P}}(m\gb)} e^{t}e_{-\gs}a_{\gs}(\mu)e^{N-t}
\nn\\&=& 
\sum_{\gs \in {{\bf P}}(m\gb)} e^{t}e_{-\gs}^\tth a_{\gs}(\mu)/{\gs(\ga)!}e^{N-t}
\nn\\&=& \sum_{\;j \geq 0}  \sum_{\gs \in {{\bf P}}(m\gb)}\left(
\begin{array}{c}
t \\
j
\end{array}
\right) j! \pd_U^{(j)}(e_{-\gs}^\tth)e^{N-j} %e^{N-mq}
a_{\gs}(\mu)/{\gs(\ga)!} 
\ey 
By condition (c) in Subsection \ref{gsi}, $\pd_U^{(j)}$ preserves the hybrid $\Z$-form on  $U_\Z(\fn^-)_e $, so 
 each term on the right side of \eqref{08x} is contained in 
$U_\Z(\fn^-)_e^{-m\gb -
N \ga }.$
 Let $q^{(j)} _{\gs}\in \Hom(U_\Z({\mathfrak n}^-)_e^{- m\gb -N\ga},\Z)$ be the function which picks out the coefficient of $e_{-\gz}^\tth$ in 
$$\left(
\begin{array}{c}
t \\
j
\end{array}
\right)  j! \pd_U^{(j)}(e_{-\gs}^\tth)e^{N-j} /{\gs(\ga)!}%e^{N-mq}
.$$
 Thus  all $q^{(j)} _{\gs}$ are simultaneously defined by 
\begin{eqnarray} %\label{b coeffs} 
\left(
\begin{array}{c}
t \\
j
\end{array}
\right) j! \pd_U^{(j)}(e_{-\gs}^\tth)e^{-j}e^{N}/{\gs(\ga)!} %e^{N-t-j} 
=\sum_{\gz \in {\overline{\bf P}}(m\gb+N\ga)}
 q^{(j)} _{\gs}( e_{-\gz}^\tth)e_{-\gz}^\tth
\nn\end{eqnarray}
and hence by 
\eqref{08x}

\by \label{107aa} e^{t} \theta_{\gb,m}(\mu)e^{N-t} =
\sum_{j,\gs,\gz} 
 q^{(j)} _{\gs}( e_{-\gz}^\tth)e_{-\gz}^\tth
a_{\gs}(\mu).
\ey


\bl \label{109b}  If  $
 q^{(j)} _{\gs}( e_{-\gz}^\tth)\neq 0,$ 
then %by Lemma \ref{1uebasis},
\begin{eqnarray} %\label{198g}
\gz(\ga)\ge N- mq.\nn
\end{eqnarray}
\el
\bpf By \eqref{107aa}  and \eqref{121nd}, since $t=p+mq$  
\by  
\sum_{j,\gs,\gz} 
 q^{(j)} _{\gs}( e_{-\gz}^\tth)e_{-\gz}^\tth
a_{\gs}(\mu) &=& e^{t} \theta_{\gb,m}(\mu)
e^{-p}e^{N-mq}\nn\\ & =& \theta_{\gc,m}(s_\ga \mu)
e^{N-mq} \in Ue^{N-mq}.\nn
\ey
This implies the result.
\epf\noi 
Hence all terms with $\gz(\ga)< N- mq$ cancel  and
using the notation of Lemma \ref{JL}  we may replace the other terms 
 $e^\tth_{ -\gz}$ by 
 $e^\tth_{- \gz'}e^{N-mq}$ where $\gz'\in \overline{\bf P}(m\gc)$.
Then from \eqref{107aa}


\by %\label{107am}
 e^{t} \theta_{\gb,m}(\mu)e^{mq-t} =
\sum_{\gz'\in \overline{\bf P}(m\gc)}\sum_{j,\gs} 
 q^{(j)} _{\gs}( e_{-\gz})e_{-\gz'}^\tth
a_{\gs}(\mu)\in U_\Z(\fn^-).
\nn\ey
\epf
\noi This means that we have no more use for the hybrid $\Z$-basis, which is  awkward to work with and use  the divided power form.   But until \eqref{09x} we have to keep track of terms that cancel. 
\bc For all $\gs,j$ as above there exist $
R _{\gs, j}\in U_\Z(\fn^-)_e^{-m\gc}$
 such that \bi\itema 
\by
%\label{xxx}
%$\quad \quad \quad \quad \quad \quad 
\left(
\begin{array}{c}
t \\
j
\end{array}
\right)  \pd_U^{j}(e_{-\gs})e^{mq-j} + R _{\gs, j} \in U_\Z(\fn^-)^{-m\gc} 
\ey
\itemb For all   
$\mu\in H_{\gb,m}$ %$$\quad \quad \quad \quad \quad \quad \quad \quad \quad \quad \quad %
\by
\label{xxx}
\sum_{\;j \geq 0}  \sum_{\gs
\in {{\bf P}}(m\gc)} R _{\gs, j}
\a _{\gs}( \mu) =0.  \ey
%IAN Do I need to also sum over j?
\ei
\ec 
\noi Let $\q^{(j)} _{\gs}\in \Hom(U_\Z({\mathfrak n}^-)^{- m\gc},\Z)$ be the function 
picks out the coefficient of $e_{- \pi}$ in 
\by
\label{A88}\left(
\begin{array}{c}
t \\
j
\end{array}
\right)  \pd_U^{j}(e_{-\gs})e^{mq-j}+R _{\gs, j} &=&\sum_{\pi
\in {{\bf P}}(m\gc)} 
\q^{(j)} _{\gs}( e_{- \pi})e_{- \pi}.  
\ey
\index{Functions $\q^{(j)} _{\gs}$} 
Then since $t=p+mq$, %we have, compare \eqref{08x} 
\by \label{09x} e^{t} \theta_{\gb,m}(\mu) 
e^{-p}
 & =& 
\sum_{\gs \in {{\bf P}}(m\gb)} e^{t}e_{-\gs} a_{\gs}(\mu)e^{-p}
\nn\\&=& \sum_{\;j \geq 0}  \sum_{\gs \in {{\bf P}}(m\gb)}\left(
\begin{array}{c}
t \\
j
\end{array}
\right)  \pd_U^{j}(e_{-\gs})e^{mq-j} %e^{N-mq}
a_{\gs}(\mu)\\&=& 
\sum_{\pi
\in {{\bf P}}(m\gc)} 
\q^{(j)} _{\gs}( e_{- \pi})e_{- \pi}a_{\gs}(\mu). \nn \ey 

\subsubsection{Proof of Theorem \ref{cot} } \label{stpre} 

We show that the coefficients $c_{ \pi}(s_\ga\cdot\mu)$  of  $\theta_{\gamma,m}$ as in \eqref{rat} are given by

\by \label{cxq}  c_{ \pi}(s_\ga\cdot\mu) %e^{t} \theta_{\gb,m}(\mu)e^{N-t} =
=\sum_{\;j \geq 0}  \sum_{\gs \in {{\bf P}}(m\gb)}
%=\sum_{j,\gs} 
 \q^{(j)} _{\gs}( e_{- \pi})
a_{\gs}(\mu).
\ey
\noi In fact this is now immediate   from \eqref{09x} and \eqref{12nd}.

\subsubsection{Related results} \label{Cons}\noi
\noi The next result is used to obtain  upper bounds on the degree  of the coefficients of \v Sapovalov elements. 
\bl \label{A7}
If  $
 \q^{(j)} _{\gs}( e_{- \pi})\neq 0,$ then
\begin{eqnarray} \label{198x}
|\pi| \leq |\gs| + mq - j.\ey
\el \bpf All terms 
$\left(
\begin{array}{c}
t \\
j
\end{array}
\right)  \pd_U^{j}(e_{-\gs})e^{mq-j} $ have 
the same degree in the standard filtration on $U({\mathfrak n}^-)_e$ including those that have poles.  
So the result holds by 
\eqref{A88}.
 \epf \noi
\bc\label{chq}\bi \itemo \itema
If $\q^{(j)} _{\gs}( e_{-\pi^0}) \neq 0,$ then ${\gs}={\gs}^0$, $j=0$ and  $\q^{(j)}_{\gs}( e_{-\pi^0}) = 1.$
\itemb If  $\a _{\gs^0} =1, $ then $c_{\pi^0}=1$.
\itemc If $\q^{(j)} _{\gs}( e_{-m\pi^\gc}) \neq 0,$ then ${\gs}={{m\gs^\gb}}$ and  $j=mq$.% and  $\q^{(j)}_{\gs}( e_{-\pi^0}) = 1.$

\ei
\ec\bpf Note that $|{\gs}|<|{\gs^0}|$ if $\gs\neq{\gs^0}$. By Lemma \ref{A7} and the hypothesis  $|\pi^0| \leq |\gs| + mq - j$. The conclusions in (a) follow easily and (b) follows from this and \eqref{cxq}. %Correct coujld explain betteer
Finally (c) holds since $e_{-m\pi^\gc} e_{-\ga}^{-mq}$ is a lowest weight vector for the adjoint action of $\fl$. The corresponding highest weight vector is $e_{-m\gs^\gb}$.% and then $j=mq$ by \eqref{198x}.
\epf \noi We also note a relation between  the functions $
 \q^{(j)}_\gs$ 
and an important  linear map   $B$. Define 
$B:
U({\mathfrak n}^-)^{- m\gb} \lra 
U({\mathfrak n}^-)^{- m\gamma}$ 
by $B(x)= x e^{mq}$. 
By \eqref{A88}\by
\label{A89}\left(
\begin{array}{c}
t \\
j
\end{array}
\right)  B(\pd_U^{j}(e_{-\gs})e^{-j}) &=&\sum_{\pi
\in {{\bf P}}(m\gc)} 
\q^{(j)} _{\gs}( e_{- \pi})e_{- \pi}.  
\ey 
%\end{document} 


%\end{document} \\ \\
%\end{document}

%\end{document}
 
\subsection{The map $B$ and the operator $\nabla^{t}_U$: noncommutative case}  \label{na1} 
\index{Maps $B_U$, $\nabla_U$: noncommutative case}
Let  $B:U_e=U_\Z(\fn^-)_e \lra U_\Z(\fn^-)_e$ be the map whose restriction to $U^{-m\gb}_e$ is given by 
\be %\label{n B} 
B(x) = xe_{-\ga}^{mq}.\nn\ee 
Let $\int$ \index{Integer valued poynomials $\int$} 
denote the ring of integer valued poynomials \index{Integer valued poynomials $\int$} in one variable This  has $\Z$-basis  
$\small 
\left(
\begin{array}{c}
T\\
n
\end{array}
\right)$, \cite{H} Lemma 26.1. We 
 define a map $\nabla=\nabla_U:U_e\ot_\Z \int\lra U_e\ot_\Z \int$ 
by
\by \label{def nabla}\nabla(x)&=&
\sum_{\;j \geq 0} 
\left(
\begin{array}{c}
T\\
j
\end{array}
\right)\pd_U^{j}(x)
 e^{-j} 
\ey
for $x\in U_e.$ Then 
for $t\in\ttk$ we obtain an operator on $U_e$ by evaluation $ \nabla^{t}_U = \nabla|_{T=t}$.  
Note that  $\nabla^{t}_U$ preserves 
 weight spaces. % we have $\nabla^{t}_U(\cO(Z)_e^{-m\gb}) \subseteq \cO(Z)_e^{-m\gc}$.
Also  %compare \eqref{b coeffs}
\by
%\label{199xx}
 B\ci \nabla^{t}_U(e_{-\gs})&=&
 \sum_{\;j \geq 0} \left(
\begin{array}{c}
t\nonumber\\
j
\end{array}
\right) B(\pd_U^{j}(e_{-\gs}) e^{-j})
\\
&=& 
\sum_{\;j \geq 0}
\left(
\begin{array}{c}
t\\
j
\end{array}
\right)   \pd_U^{j}(e_{-\gs}) e^{mq-j}.\nn
\ey
Combined with \eqref{1cow}, this gives
\be \label{fxx} 
\Phi_t(e_{-\gs})e^{mq} = B\ci \nabla^{t}_U(e_{-\gs})
.\ee

\subsection{Reordering PBW bases}\label{zzw} To prove the results on coefficients as soon as possible, we  quote a result out of sequence. 
\bc  \label{3W1}  For any rogue partition 
$\pi$ of $m\gc$ and any order on positive roots, the coefficient of $e_{-\pi}$ in $\theta_{\gc,m}$ is zero.
\ec \bpf This follows 
from 
Corollary \ref{3W} (c)
since if $a\in S(\fh)$ is nonzero, then $\gk(a)$ is nonzero.
\epf
\noi Without loss of generality we can work over $\ttk$.
  Consider  two orders on $\Delta^+$, and for $\pi \in
{\bf P}(m\gc)$, set $e^\ttk _{-\pi} = \prod_{\alpha \in \Delta^+} e^{\pi
(\alpha)}_{- \alpha}$, and $\overline{e}^\ttk _{-\pi} = \prod_{\alpha \in
\Delta^+} e^{\pi (\alpha)}_{- \alpha}$, the products being taken with
respect to the given orders. Then we have two bases $\{
e^\ttk _{-\pi}|\pi \in {{\bf P}}(m\gc)\}$  and $
\{\overline{e}^\ttk_{-\pi}|\pi \in {{\bf P}}(m\gc)\}$  for $U(\fn^-)^{-m\gc}$.  However  to write $\theta_{\gc,m}$ we do not need rogue partitions. So define
 $${{\bf Q}}(m\gc)=\{\pi \in {{\bf Q}}(m\gc)| \pi 
\mbox{ is not a rogue partition of } m\gc
\}.$$ We have linearly independent sets ${\bf B_1} = \{
e^\ttk _{-\pi}|\pi \in {{\bf Q}}(m\gc)\}$  and ${\bf B_2} = \{\overline{e}^\ttk_{-\pi}|\pi \in {{\bf Q}}(m\gc)\}.$  
Fix a total order on the set ${{\bf Q}}(m\gc)$ such that if $\pi, \gs \in {{\bf Q}}(m\gc)$ and $|\pi| > |\gs|$ then $\pi$ precedes $\gs,$ and use this order on partitions to induce orders on the bases ${\bf B_1},{\bf B_2}$.
\bl \label{der1}   We have $\span \;{\bf B_1} =\span \;{\bf B_2}$
and the change of basis matrix $C=(c_{\pi,\zeta})$ from the basis ${\bf B_2}$ to ${\bf B_1}$ 
is an upper triangular with all diagonal entries equal to $\pm1.$ \el %
\noindent \bpf This follows since %from Corollary \ref{3W1} and the fact that  
the associated graded ring arising from the standard filtration on
$U=U({\mathfrak n}^-)$ is supercommutative.
\epf

\bexa \label{ex1}{\rm Suppose that $\fg = \osp(5,2)$ or $\osp(4,2)$.  In the former case, $\fg$ has a basis of simple roots consisting of the roots $\ga_1= \epsilon_{1} - \epsilon_{2}$, $\ga_2=\epsilon_{2} - \delta $ and $\ga_3=\delta$.  In the latter case $\delta$ is replaced by {$\ga_3=2\delta$}.  For the corresponding Dynkin-Kac diagrams see \cite{M101} Tables B.5 and
B.7. Consider the partitions $\pi, \gs $ of  $\gc =
\epsilon_{1} + \epsilon_{2}$ taking nonzero values 
$ \pi(\epsilon_{1} - \gd) =\pi(\epsilon_{2} + \delta)=1$  
and $\gs( \epsilon_{1} + \epsilon_{2})=1.$  For different orders on the positive roots we have
% $e_{-\pi} =e_{\gd-\epsilon_{1} } e_{\epsilon_{2} - \delta}$,   $\overline{e}_{-\pi} = e_{\epsilon_{2} - \delta}e_{\gd-\epsilon_{1} } $
 $e_{-\pi}^\ttk=e_{\gd-\epsilon_{1} } e_{-\epsilon_{2} - \delta}$,   
$\overline{e}^\ttk_{-\pi} = e_{-\epsilon_{2} - \delta}e_{\gd-\epsilon_{1} } $
and $e_\gs^\ttk =\overline{e}^\ttk_\gs =e_{ -\epsilon_{1} - \epsilon_{2}}$.  We have,  %can assume that 
$e_{-\pi}^\ttk + \overline{e}_{-\pi} = c\overline{e}_{-\gs}$ for a nonzero constant $c$, or in matrix form
$$\left[\begin{array}{c} e^\ttk_{-\pi}\\ {e^\ttk}_\gs \end{array}\right] =\left[\begin{array}{cc} -1&c\\ 0&1 \end{array}\right]
\left[\begin{array}{c}  \overline{e}^\ttk_{-\pi}\\ \overline{e}^\ttk_{-\gs}\end{array}\right].$$
}
\eexa\noi 
Consider a basis 
$\cB = \{{b_\pi}|\pi \in {{\bf Q}}(m\gc)\}$  indexed by partitions.  
Fix a total order on the set ${{\bf Q}}(m\gc)$ such that if $\pi, \gs \in {{\bf Q}}(m\gc)$ and $|\pi| > |\gs|$ then $\pi$ precedes $\gs,$ and use this order on partitions to induce orders on the basis ${\cB}$.  If 
$x \in U(\fb^-)^{-m\gc}$ we can write $x$ uniquely in the form

\be %\label{owl1} 
x \; = \; \sum_{\pi \in {{\bf Q}}(m\gc)}
b_\pi h_\pi \nn.\ee 
We list various desirable conditions on the basis $\cB.$
   \bi \itema
$\deg  h_\pi \le m\hgt \gc - |\pi|$  
for all $\pi \in {{\bf Q}}(m\gc),$
\itemb As a polynomial function of
$\gl$,  
$h_{m\pi^{\gamma}}$ has leading term  $c\prod_{\gs
\in N(w^{-1})}(\gl,\gs)^{mq(w,\gs)}$ for a nonzero constant $c$.
%\itemc If $\pi \in {{\bf R}}(m\gc),$ then $h_{\pi} =0.$ 
 \ei  


\bl \label{derx}
For $x \in U({\mathfrak n}^-)^{-m\gc}
\otimes S({\mathfrak h}),$ and using the bases ${\bf B_1}$ and  ${\bf B_2}$ write
\be %\label{1owl} 
x \; = \; \sum_{\gz \in {{\bf Q}}(m\gc)}
\overline{e}^\ttk_{-\gz} f_\gz \; = \; \sum_{\pi \in {{\bf Q}}(m\gc)}
{e}_{-\pi}^\ttk g_\pi.\nn\ee  Then  the basis ${\bf B_1}$ satisfies conditions $(a)-(b)$ iff the basis  ${\bf B_2}$ 
satisfies the corresponding condition. 
\el
 \bpf  Assume ${\bf B_2}$ satisfies conditions $(a)-(b).$ Using the matrix $C$ from Lemma \ref{der1} we can write \[{e}^\ttk_{-\pi} = 
\sum_{\zeta} c_{\pi,\zeta}\overline{e}^\ttk_{-\gz}\] 
 Then %Thus (\ref{1owl}) holds with
$$f_\zeta = \sum_{\pi \in {{\bf Q}}(m\gc)}
c_{\pi,\zeta}g_{\pi}.$$
 It follows that $f_\zeta$ is a linear combination of polynomials of
degree less than or equal to $m\hgt \gc - |\gz|.$ 
Thus ${\bf B_1}$ 
 satisfies (a). Also the partition 
 $m\pi^{\gc}$ occurs last in both orderings, and ${e}^\ttk_{-m\pi^{\gc}} =e_{-\gc}^m=\overline{e}^\ttk_{-m\pi^{\gc}}$ occur with the same coefficient. Thus ${\bf B_1}$ 
 satisfies (b).\epf

\subsection{Proof of Theorem \ref{1sag}} \label{dt}  
By Lemma \ref{derx}, it is enough to prove  Theorem \ref{1sag} when we fix an $\ga$-order on $\Gd^+$.
\noi  Write \be %\label{1wuga}
 w =
s_{\ga}u, \quad   \gb = u{\gb_0}, \quad \gamma = w{\gb_0} = s_{\alpha} \gb,
\nn\ee
 with $\alpha \in \Pi_{\rm nonisotropic}$ and $\ell(w) =\ell(u) + 1.$
Suppose $\nu \in \Lambda$ and set  
\be %\label{1wuga1}
\mu =
u \cdot\nu, \quad \lambda = w \cdot\nu
= s_{\alpha} \mu.
 \ee
 By induction, $$\deg  a_{\gs}(\mu) \le m\hgt \gb - |\gs|.$$ By Theorem   \ref{cot}, 
\be \label{sas}\mbox{deg} \; c_{\pi}(\gl)  \leq  \max\{j + \; \mbox{deg} \;
a_{\gs}(\mu)\;|\; \q^{(j)} _{\gs}( e_{- \pi})
\neq 0\}.\ee
  Now if $\q^{(j)} _{\gs}( e_{- \pi})
\neq 0$ then by (\ref{198x}) %with $N=mq,$
$j \leq |\gs| +mq -|\pi|.$  %
We eliminate $j$ by combining  this inequality with \eqref{sas}, giving the first inequality  below. The second holds by induction and the last equality by  \eqref{121d}
\begin{eqnarray} %\label{1cat}
\mbox{deg} \;\; c_{\pi}(\gl) \nonumber &\leq&
\mbox{deg} \;\; a_{\gs}(\mu)+ |\gs| + mq - |\pi|
\\ &\le& m \hgt \gb - |\pi| +mq\nonumber\\
&=&  m\hgt \gc - |\pi|.
%\ee
\nn\end{eqnarray}
\subsubsection{Leading term} \label{tdt} First note that by \eqref{iv} we have 
\be %\label{s13}
q(w,\ga) = (\gc, \alpha^\vee) = q\nn\ee and 
\be %\label{s12} 
q(u,\gt) = (u\gb, \gt^\vee) =  (w\gb, s_\ga \gt^\vee)= q(w,s_\ga \gt) .\nn\ee  \noi  
It follows from Theorem \ref{cot} amd  Corollary  \ref{chq} that, modulo terms of lower
degree
\be %\label{1Nw1} 
 c_{m\pi^{\gamma}}(\gl) \equiv \left( \begin{array}{c}
t\\
mq
\end{array}   \right)
a_{m\gs^{\gb}}(\mu) .\nn\ee
 By induction %$a_{m\gs^{\gb}}(\mu)$
 $a_{m\gs^{\gb}}(\mu)$ has
highest term $c'\prod_{\gt \in N(u^{-1})}(\mu,\gt)^{mq(u,\tau)}$ as a polynomial in $\mu$, for
a nonzero constant $c'$.   
Thus  as a polynomial in $\gl=s_{\alpha} \mu$, $a_{m\gs^{\gb}}(\mu)$ has
highest term 
\by %\label{1Nw} 
c'\prod_{\gt \in N(u^{-1})}(\gl,s_\ga\gt)^{mq(u,\tau)} 
&=&c'\prod_{\gs \in s_\ga N(w^{-1}), \gs\neq \ga}(\gl,\gs)^{mq(w,\gs)},\nn\ey 
On the other hand $t=\pm (\gl,\ga) + C$, with some constant $C$ independent of $\mu, \gl$.  Thus
\be %\label{1Nw2}  %c_{m\pi^{\gamma}}(\gl) \equiv 
\left( \begin{array}{c}
t\\
mq
\end{array}   \right)\equiv t^{mq}\equiv (\gl,\ga)^{mq}
\nn \ee
modulo terms of lower degree and up to a nonzero constant multiple. The
claim about the leading term of $c_{m\pi^{\gamma}}(\gl) $ in Theorem \ref{1sag} follows from the preceding
equations and \eqref{Nw}.% \eqref{s13}, \eqref{s12}, (\ref{1Nw1}) and \eqref{1Nw}.
%For $G_2$ with $\ga$ short $\gb$ long, the highest short root is $2\ga+\gb=s_\ga(\ga+\gb)=  s_\ga s_\gb\ga$, take $w=s_\ga s_\gb$and note $t$ in the last step is $\ga+\gb$ and $N(w^{-1}) =\{ \ga, \ga +\gb\}$

\subsection{Formula for coefficients}\label{cfw}
Here we work with an $\ga$-order on positive roots. %IAN this seems to be a restatement of Theorem \ref{cot}.

\bt \label{cgw} If $t = mq+p$, as in  \eqref{121d}, 
then the  Šapovalov elements $\theta_{\gamma,m}$ and $\theta_{\gb,m}$ are related after evaluation on $\mu\in{H}_{\gb, m}$
by 
\by 
\label{177} \sum_{\pi \in {{\bf P}}(m\gamma)}  (1\ot P_\ga) e_{-\pi} c_\pi
&=& \sum_{\gs \in {{\bf P}}(m\gb)} ((B_U\ci\nabla_U^t)\ot 1)e_{-\gs} a_\gs. 
\ey
\et
\bpf After evaluation, % both sides on $\mu\in{H}_{\gb, m}.$
by \eqref{rat} (with 
$\gl=s_\ga\cdot\mu$) 
%\eqref{rat} 
and \eqref{12md} the left side  becomes
$\Phi_t(\theta_{\gb,m}(\mu)  )e^{mq}.$ By \eqref{1x7}  the right side becomes $B\ci\nabla^{t}_U( \theta_{\gb,m}(\mu))$. 
These are equal by \eqref{fxx}.
\epf\noi 
We can also write 
\eqref{177} in the form 

\by 
\label{17x} \sum_{\pi \in {{\bf P}}(m\gamma)}  (1\ot P_\ga) \go(x_{\pi}) c'_\pi
&=& \sum_{\gs \in {{\bf P}}(m\gb)} ((B_U\ci\nabla_U^t)\ot 1)\go(x_{\gs}) a'_\gs. 
\ey where $\gk(c'_\pi)= \gk(c_\pi)$ and $\gk(a'_\gs) =\gk(a_\gs) $,
\subsection{Uniqueness of \v Sapovalov elements} \label{1s.1}

 The \v Sapovalov determinant $\det F_{\eta}$, the Jantzen filtration $\{M_{i}(\lambda) \}$ and sum formula are defined as usual, \cite{M101} Chapter 10.  
 %They can be used to  prove a uniqueness result for \v Sapovalov elements.  
Proofs are mostly omitted as they are %essentially 
the same as % very similar to those 
in \cite{Gk},   \cite{KK} or  \cite{M101}.   Recall the definitions of
$A(\lambda), B(\lambda)$ from equations (\ref{1tar}), (\ref{Bdef}).
%\vspace{0.5cm}\\
%{\underline{\bf Lemma.}}
\begin{theorem} \label{shapdet}
Up to a nonzero constant factor, $\det F_{\eta} = D_{1}D_{2}D_{3},$
where
\begin{eqnarray*}
D_{1} & = & \prod_{\alpha \in \overline{\Delta}^+_{0}} \;\;
\prod^\infty_{r=1}
            (h_{\alpha} + (\rho, \alpha) - r(\alpha,\alpha)/2)^
            {{\bf p}(\eta - r \alpha)}, \\
D_{2} & = & \prod_{\alpha \in
\Delta^+_{1}\backslash\overline{\Delta}^+_{1}}\;\;
            \prod^{\infty}_{\stackrel{\scriptstyle{r=1}}{\scriptstyle{r \; {\rm odd}}} }
            (h_{\alpha} + (\rho, \alpha) - r(\alpha,\alpha)/2)^
            {{\bf p}(\eta - r \alpha)}, \\
D_{3} & = & \prod_{\alpha \in \overline{\Delta}^{+}_{1}} \;\;
(h_{\alpha} +
             (\rho, \alpha))^{{\bf p}_{\alpha}(\eta - \alpha)}.
\end{eqnarray*}
\end{theorem}
%\newline
\noi This is proved by first finding the leading term and then using representation theory to analyze $M(\gl)$ in the most generic cases.
Suppose  $\alpha$ is a positive nonisotropic root of $\mathfrak{g}$, and $r \in
\mathbb{Q} ^+$ and $ (\lambda + \rho, \alpha^{\vee} )  =  r $.
Then either $M(\lambda)$ is simple, or
$M(\lambda)$ has a submodule $U$ which is isomorphic to
a finite direct sum of modules isomorphic to
$M(s_\alpha \cdot \lambda)$ and $M(\lambda)/U$
is simple, \cite{KK} Lemma 3.3 (b), 
\cite{M101} Lemma 10.2.4.  However using Lemma \ref{xyz} it follows that if $M(\lambda)$ not simple, then $U\cong M(s_\alpha \cdot \lambda)$. Later in the proof it is shown that 
$\ga\in A(\lambda)$.
%\\ \\\noi A remarkable feature of the formula in Theorem \ref{shapdet} that it is stable under odd reflections.  Note that odd reflections contribute to the factor $D_3$ and if $\ga$ is a simple isotropic root then $D_3$ is unchanged if $\ga$ is replaced by $-\ga$.
\\ \\
We now state the Jantzen sum formula for the Verma
modules $M(\lambda)$.  
\begin{theorem} \label{Jansum}
For all $\lambda \in \mathfrak{h}^*$
\[ \sum_{i > 0} \ch M_{i}(\lambda) = \sum_{\alpha \in A(\lambda)}
\ch M(s_\alpha \cdot  \lambda) + \sum_{\alpha \in
B(\lambda)} \epsilon^{\lambda -\alpha}p_\ga.\]
\end{theorem}
\bt \label{17612} Suppose $\gth_1, \gth_2$ are \v Sapovalov elements for the pair $(\gc,m)$.  Set $H = H_{\gc, m}.$ Then
for all $\gl \in H$ we have $\gth_1 v_\gl =\gth_2 v_\gl$. 
%\itemb $\gth_1- \gth_2\in U({\mathfrak g}){\mathfrak n}^+  + U({\mathfrak g}){\mathcal I}(H).$ \ei
\et \bpf
Let $\{{M}_{i}(\lambda)\}_{i\ge 1}$ be the Jantzen filtration on 
${M}(\lambda)$. Then ${M}_{1}(\lambda)$ is the unique maximal submodule of ${M}(\lambda)$.  
Set
\[\gL = \{\gl \in H| A(\gl) = \{\gc\}, \;B(\gl) = \emptyset \},\]
if $\gc$ is non-isotropic, and
\[\gL = \{\gl \in H| B(\gl) = \{\gc\}, \;A(\gl) = \emptyset \},\]
if $\gc$ is isotropic.  If $\gl \in \gL$
it follows from Theorem \ref{Jansum}, that
$ M_1(\gl)^{\gl-m\gc}$ is one-dimensional.
Because ${M}_{1}(\lambda)$ is the unique maximal submodule of
${M}(\lambda)$, $\gth_1 v_\gl$ and $\gth_2 v_\gl$ are proportional. Then from the requirement that $e_{-\pi^0}$ occurs with coefficient 1 in
a \v{S}apovalov element we have $\gth_1 v_\gl =\gth_2 v_\gl$.  Since $\gL$ is Zariski dense in $H$, the result 
follows.
\epf
\noi 
A simple example shows that uniqueness of \v Sapovalov elements fails if $\fg(A,\gt)$ is replaced by a subfactor.  Consider the matrix
$$A= \left[\begin{array}{ccc} 2&-1&0\\-1& 0&1\\0&1&-2 \end{array}\right]$$
If $\gt =\{2\}$, then $\fg(A,\gt)\cong \fgl(2|2). $ The subfactor  
$\fg=\fpsl(2|2)$ has simple roots $\ga, \gb, \gc$, where $\ga+ 2\gb+ \gc=0.$ 
If we use the Borel subalgebra with simple root vectors $e_{13}, e_{34}, e_{42},$ then $e_{13}$ and $e_{42}$ have the same weight $\ga+\gb $, so if $(\gl,\ga+\gb)=0$ we have $$\dim \Hom_{U(\fg)}(M(\gl-(\ga+\gb)), M(\gl))=2.$$ 
\noi %\end{document}

%\end{document}- 
\section{Commutative version} \label{g1}

\subsection{The operator $\nabla_S$: commutative case}  \label{na2}
Recall the derivations $ \pd_S^{\pm}$ of  $S(\fn^-)$ defined in Subsection \ref{taa} and 
define derivations  $d_{\pm}$ of  $S(\fn^-)_e$  by 
\be\label{A2} d_{\pm} = x_{\ga}^{\pm 1} \pd_S^{\pm}.\nn\ee 
Recall that $x_{\ga}$ is a copy of $e_{-\ga}$ in $S(\fn^-)$.
\index{Derivations $d_\pm$} 
Let $\Div$ be the ring of divided powers of $T$. \index{Ring  of divided powers $\Div$}
This has $\Z$-basis  
 $T^{(n)}= \frac{T^n}{n!}$
  for $n\ge 0$.
Define a map $\nabla_S:S(\fn^-)_e\ot \Div \lra S(\fn^-)_e\ot \Div$ by
\by \label{Sdef nabla}\nabla_S(x) = 
(\exp T d_-)(x).
\ey
for $x\in S(\fn^-)_e.$ Then 
for $t\in\ttk$ set  $ \nabla^{t}_S = \nabla_S|_{T=t}=\exp t d_-$.    
\index{Maps $B_S$ and $\nabla_S$: commutative case}  It follows from (c) in Subsection \ref{gsi}, that $\exp   t d_\pm$ preserves $S_\Z(\fn^-)_e$.
\bl\label{dra}\bi \itema $\nabla_S$ and $\nabla^t_S$ are graded algebra maps.
\itemb We have $\nabla^{t}_S\ci \nabla^{s}_S = \nabla^{t+s}_S$ for all $s,t\in \ttk$.
\itemc $( \nabla_U\ci\go)(u) \equiv (\go \ci\nabla_S)(u)$ modulo lower degrees in $T$.%the filtration on $S(\fh)\ot \ttk[T].$
\itemd For $x$ as in \eqref{nfn} we have $( \gk\ci\nabla_U\ci\go(x)) \equiv (\gk\ci\go \ci\nabla_S)(x)$ mod lower degrees in $T$.
\ei
\el \bpf  First (a) and (b) holds for the exponential of any locally nilpotent derivation.  
%For (c) if $x$ as in \eqref{nfn},  
Since $\go\pd^j_S(u) = \pd^j_U\go(u)$ by Lemma \ref{npr} 
\by \label{dea}\nabla_U\ci
\go(u)&=&
\sum_{\;j \geq 0}  \left(
\begin{array}{c}
T\\
j
\end{array}
\right)
\pd_U^{j } (\go(u)
 e^{-j})
\nn\\ 
&=& \sum_{\;j \geq 0} \left(
\begin{array}{c}
T\\
j
\end{array}
\right)
\go(\pd_S^{j}(u))
 e^{-j}.\nn
\ey
 A similar calculation shows that 
\by \label{bea} \go\ci\nabla_S (u)%&=&\sum_{\u}\nabla^{t}_S(u_\u)\gk(a_\u)\nn\\ 
&=& \go\ci\sum_{\;j \geq 0} \frac{T^j}{j!}\pd_S^{j}(u)
 e^{-j} \nn\\
&=&
 \sum_{\;j \geq 0} \frac{T^j}{j!}\go\pd_S^{j}(u)
 e^{-j}.
\ey 
This proves (c) and  (d)  is an immediate consequence.
\epf \noi

%$fff=x_\gc^{(m)} = x_\gc^m/m! $%

\subsection{ \v Sapovalov subalgebras}  \label{ve} 
 Note that  $S(\fn^-)_e$  has a weight space decomposition.  Fix  $\gt\in \{\gb,\gc\}$. 
 The $\gt$-Šapovalov algebra is the subalgebra of $S(\fn^-)_e$ given by  %depends on an $\ga$ order on positive roots. First we define the 
$$
\mbox{Š}_\gt
= \bigoplus_{m\ge 0}
\mbox{Š}_\gt(m), \quad 
\mbox{ where }\mbox{Š}_\gt(m)= (S(\fn^-)_e)^{-m\gt}.$$ 
 Note that Š$_\gt$ is a graded algebra, but if $\gt$ is not a simple root, the grading is not compatible with the usual grading on $S(\fn^-)$.   What makes up for this is that  all the algebras are graded by the root lattice and this grading behaves well. Let $\mbox{Š}^0_{\gt}$ be the subalgebra of $\mbox{Š}_{\gt}$ generated by $\mbox{Š}_{\gt}(1)$. \index{$\gt$-Šapovalov algebra!integral form} The integral forms of  $\mbox{Š}_{\gt}$ and  $\mbox{Š}^0_{\gt}$
are defined by $\mbox{Š}_{\gt,\Z} = \mbox{Š}_{\gt}\cap S_\Z(\fr)[e^{\pm 1}]$
and  $\mbox{Š}^0_{\gt,\Z}= \mbox{Š}^0_{\gt}\cap S_\Z(\fr)[e^{\pm 1}]$
respectively. 
 Both derivations  $d_\pm$
preserve weight spaces and the integral forms $\mbox{Š}_{\gb,\Z}$. 
%\index{Integral forms! $\mbox{Š}_{\gb,\Z} , \mbox{Š}_{\gc,\Z}$}
\\ \\
From now on we restrict  $\nabla_S=\exp T d_-$ to Š$_\gb[T]$.  Obviously $\nabla_S$ is  an algebra automorphism  
of Š$_\gb[T]$.  Also there is  an isomorphism  of graded algebras 
\be\label{BSD} B_S:\mbox{Š}_\gb[T]\lra \mbox{Š}_\gc[T]\ee which is given
on Š$_\gb(m)$ by $x\lra xe^{mq}$. %Note  the multiplication here takes place in $S(\fn^-)_e$   In the next section we study the functors on module categories

%$\cO_\Z(H_{\gb,m}), \cO_\Z(H_{\gc,m})$

\noi \subsection{Module categories} \label{PA.101}

\subsubsection{Induction and restriction} \label{td2} 
Let $f¨:R\lra A$ be a honomorphism of rings. Denote their module categories by $R$-mod and $A$-mod respectively. For $N\in R$-mod denote the induced $A$ module $\inr_R^A\; N$. 
If $M\in A$-mod we obtain an $R$-module   $\rer_R^A \;M$ and we  have a natural isomorphism 
\be \label{GF11}\Hom_A(\inr_R^A\; N, M) 
\cong  \Hom_R(N, \rer_R^A \;M).\ee
We denote the functors $\inr_R^A\; ,\rer_R^A $ respectively by $f^*$ and $f_*$. 
Taking $M=\inr_R^A\; N$ and the identity map on the left side of \eqref{GF11},  we obtain an adjunction isomorphism of $R$-modules  
$N\lra f_* f^* N.$ Explicitly $N \lra 1\ot N  \subseteq A\ot_R N.$ Taking $R=N$ we have $R \lra 1\ot R  \subseteq A\ot_R R \cong A.$
\subsubsection{Definition of $B_S$} \label{td1} 
We now apply this to the rings
$R=\mbox{Š}_{\gb,\Z}, S= \mbox{Š}_{\gc,\Z}$, $A=S_\Z(\fn^-)_e$ and the  inclusions $i:R\lra A$ and $j:S \lra  A$.
By adjunction we have $R \lra 1\ot\oplus_{m\ge 0}R(m)$. 
Since $i^*R$ is an $A$-module it is legitimate to form the sum
$\oplus_{m\ge 0}e^{mq}\ot R(m)$
which maps to $\oplus_{m\ge 0}e^{mq}
 R(m)
\subseteq (S(\fn^-)_e)^{-m\gc}=$Š$_\gc(m)$. It is easy to see the inclusion is an equality.  However we have suppressed the restriction  $j_*$ (as often happens).  To do this more formally, let $N= \op N(k) $ be a graded $R$-module and define a map $B_k: 1\ot N(k) \lra e^{kq}\ot N \lra e^{kq} N$.  Then $B_*: {\bf gr} R\lra {\bf gr} S$ is the map  
which restrcited to degree $k$ is given by 
\be \label{ma4}
N(k)\stackrel{i_* i^*}{\lra} 1\ot N(k)\stackrel{B_k}{\lra} e^{kq} \ot N(k) \stackrel{j_* }{\lra}  e^{kq}  N(k).\ee
In particular, taking $N=R$ gives the map $B_S: R\lra S$.
\\ \\
The functor  $(\nabla^t_S)_*$ is much easier to describe. Indeed $(\nabla^t_S)_*N$ is just the $R$ module 
$N$ twisted by the automorphism $(\nabla^t_S)_*$. Thus $(r,n)\lra \nabla^t_S(r)n$.
%\subsubsection{The Serre subcategory of torsion modules} \label{tdt3} 
\subsubsection{Quotient categories} %\label{td2} 
Let 
 $N=\bop_{n\ge n_0} N(n)$ be a graded $R$-module. We say that $N$ is bounded if $N(n)=0$ for $n\gg 0$ and {\it torsion} if it is a direct limit of bounded modules. Let ${\bf gr} R$ be the category of graded  $R$-modules 
and $\tors$ the full subcategory of torsion modules. 
\index{Category of graded $R$-modules ${\bf gr} R$}
\index{Category of graded $R$-modules ${\bf gr} R$!mod torsion ${\bf gr} R/\tors$}
We form the quotient category ${\bf gr} R/\tors$. The irrelevant ideal of $R$ is  $R_+=\bop_{n> 0} R(n).$
\bt \label{HF9} If $R=\mbox{Š}^0_{\gb,\Z}, S= \mbox{Š}^0_{\gc,\Z}$, the maps $\nabla^t_S$ and $ B_S$ induce equivalences of categories 

\be \label{LF9}{\bf gr} R  \stackrel{(\nabla^t_S)_*}{\lra} {\bf gr} R  \stackrel{B_*}{\lra}{\bf gr} S\ee  %\index{Category of graded $A$-modules $Gr\,A$}
and 
\be \label{LF10}{\bf gr} R/\tors  \stackrel{(\nabla^t_S)_*}{\lra} {\bf gr} R/\tors  \stackrel{B_*}{\lra}{\bf gr} S/\tors\ee
\et
%\index{Category of graded $A$-modules $Gr\,A$!mod truncation ${\bf Gr}\,A$}
\noi 
If $\gb$ is  isotropic then $\mbox{Š}^0_\gb$ is finite dimensional and local. Thus the irrelevant  ideal is torsion and ${\bf gr} \;\mbox{Š}^0_\gb/\tors  ={\bf gr} \;\mbox{Š}^0_\gb(0)/\tors  $.
\subsection{Proof of Theorem \ref{783}}  \label{GF2}

If $I_m=S_\Z (\fh) \cap I( H_{\gb,m})$,
 we define $\cO_\Z(H_{\gb,m})=S_\Z (\fh) /I_m$. \index{Integral forms!$\cO_\Z(H_{\gb,m}), \cO_\Z(H_{\gc,m})$}
This  is a $\Z$-form of $\cO(H_{\gb,m})$. 
%Let $R=\mbox{Š}_{\gb,\Z}, S= \mbox{Š}_{\gc,\Z}$ as in the previous Subsection. %Define \be \label{vat}{ H}_{\gb, m} = \{ \lambda \in  {\mathfrak h}^*|(\lambda, \gb) = m(\gb, \gb)/2  \}.\ee
% degree $m$ part
 for all $m\ge 0$.  
For  $\gt, \gs\in \{\gb,\gc\}$ we define a \index{Algebras $R(\gt,\gs)$!
$R_\Z(\gt,\gs)$} graded ring
$R_\Z(\gt,\gs)= \mbox{Š}^0_{\gt,\Z}\ot_\Z \cO_\Z(H_{\gs,0})$
 which has for all $m\ge 0$, degree $m$ part 
 % $ R_Z(\gc,\gc)$, 
$$ R_\Z(\gt,\gs)(m)= \mbox{Š}^0_{\gt,\Z}(m)\ot_\Z \cO_\Z(H_{\gs,m})$$
We use Lemma \ref{td5} to multiply the second factor. 
We have  homomorphisms  
\be \label{zra1} 
(1\ot P_\ga):  R_\Z(\gc,\gc)\lra  R_\Z(\gc,\gb)\ee and 
\be \label{zra2}  (B_S\ci\nabla_S^t\ot 1): R_\Z(\gb,\gb) \lra  R_\Z(\gc,\gb).\ee

\bt \label{zra} Assume $\gb$ is a real root. If $t = mq+p$, as in  
\eqref{121d}, 
then the leading terms of the coefficients of the Šapovalov elements $\theta_{\gamma,m}=\sum_{\pi }  x_{\pi} c_\pi$ and $\theta_{\gb,m}
=\sum_{\gs }  x_{\gs} a_\gs$ are related by
\by 
\label{1789} \sum_{\pi\in {{\bf P}}(m\gc)}  (1\ot P_\ga) x_{\pi}\gk( c_\pi)
&=& \sum_{\gs \in {{\bf P}}(m\gb)} ((B_S\ci\nabla_S^t)\ot 1)x_{\gs} \gk( a_\gs). 
\ey 
\et \bpf 
Apply the left side  of Lemma \ref{dra} (d) to the left side of the  noncommutative version, Equation   \eqref{17x}  and apply the right side  to the right side of 
 \eqref{17x}.
\epf
\noi 
By \eqref{zra1}
 and \eqref{zra2} we have an isomorphism 
\be \label{2W4}f: R_\Z(\gb,\gb)
\lra R_\Z(\gc,\gc).\ee where
$$f:=(1\ot P_\ga)(B_S\ci\nabla_S^t\ot 1) = (B_S\ci\nabla_S^t\ot 1)(1\ot P_\ga).$$
Then by Theorem \ref{zra} 
\be \label{za1} 
f( \gk(\theta_{\gb,m})) =\gk(\theta_{\gamma,m}).\ee
Taking into account the gymnastics necessary to define $B_S$, 
we note that $f$ factors through $F$, where  
\be \label{2W9}F=(1\ot P_\ga)(B_S\ci\nabla_S^t\ot 1):
R_\Z(\gb,\gb)\lra S_\Z(\fn^-)_e \ot \cO_\Z(H_{\gc,0}).\ee
By induction and \eqref{za1}, we have 
\bc  \label{3W} \bi \itema If $\gc$ is isotropic then  $\theta_{\gamma}^2 =0.$
\itemb If $\gc$ is nonisotropic and $m$  satisfies condition (b) before  Theorem \ref{1sag} then $\gk(\theta_{\gamma,m}) = \gk(\theta_{\gamma,1})^m$.
\itemc For any rogue partition 
$\pi$ of $m\gc$ the coefficient of $e_{-\pi}$ in $\theta_{\gc,m}$ is zero.
\ei
\ec \bpf Since $f$ is an algebra map, this follows since $\gk( c_\pi) =0$ in implies  $ c_\pi =0$. \epf
\noi See \cite{Md} for a different approach.
\bc \label{3W2} We have equivalences of categories and a commutative diagram. 
$$
\xymatrix{{\bf gr} R(\gc,\gc)/\tors   \ar@/_1pc/[rr]_{{(B_S\ci\nabla_S^t\ot 1)_{*}}} %\ar@/^/[drr]^x
\ar@{>}^{f_{*}}[r]
&{\bf gr} R(\gb,\gb)/\tors  \ar@{>}^{{(1\ot P_\ga)}_{*}}[r]&{\bf gr}R(\gb,\gc)/\tors .}$$ 
\ec

%\mbox{Š}_{\gc}$ generated by the degree 1 part $\mbox{Š}_{\gc}(1)$. 
\subsection{Finite generation}  
\label{2W} \noi The next result might appear redundant since we mainly work with small Šapovalov algebras, which are generated by their the space of elements of degree $\le 1$. 
However we could expect to obtain other results by applying convex geometry to
Šapovalov algebras.  For example they should be Cohen-Macaulay rings, \cite{BH} Chapter 6.
Assume that $\fg $ is finite dimensional. 
\bl \label{2W1} The algebras Š$_\gc$  and Š$_\gb$  are finitely generated. %(and hopefully normal and Cohen-Macaulay.)
\el
\bpf We prove the result for Š$_\gc$.  The first step is to show that Š$_\gc$ is isomorphic to a certain semigroup algebra $\ttk S$.  Let 
$S(m) =\{(k,\pi)\in {\Z\ti
\bf{\overline{P}}}(m\gc-k\ga)\},$ using the notation 
from Lemma \ref{1uebasis}.
Then $S= \op_{m \ge} S(m)$ is an additive semigroup, such that Š$_\gc=\ttk S$ the semigroup algebra.
 Let $\cA=\{\ga_1,\ldots \ga_R\}$ be the set of positive roots such that for each $i$ there exists an integer $k$ we have $\ga_i < k\gc$. This means that there is a partition of  $k\gc$ involving $\ga_i.$
Let $M=\sum_{i=1}^R \Z v_i$ be a free abelian group of rank $R$. 
 Assume $\ga=\ga_R$ and set 
$$H= \R v_R +\sum_{i=1}^{R-1} \R_{\ge 0} v_i \subset \R^R =M\ot \R.$$
For $t\in  \R_{\ge 0}$,  set % consider the convex set  $F:=\sum F_m$ where  
$$F_{t}=\{\sum a_i v_i\in H| \sum a_i \ga_i =t\gc\}.$$ 
and  consider the convex set  $F:=\bcu_{t\in  \R_{\ge 0}} F_t.$  
Let 
$S(m) =\{(k,\pi)\in {\Z\ti
\bf{\overline{P}}}(m\gc-k\ga)\}$. Then $S= \op_{m \ge 0} S(m)$ is an additive semigroup, such that Š$_\gc=\ttk S$ the semigroup algebra.  
For $m\in \Z^+$ and  a pair $(k,\pi)\in {\Z\ti
\bf{\overline{P}}}(m\gc-k\ga)$ as in Lemma \ref{1uebasis} we associate the point 
$v_{k,\pi} = k v_R + \sum_{i=1}^{R-1} \pi(\ga_i) v_i \in H\cap M.$ 
This gives a surjective map from the basis of  $S_\gc(m)$ given in \eqref{10b} to $F_m\cap M$ which extends to an homomorphism of semigroups $\phi:S\lra F\cap M.$  \\ 
\\
We show $F\cap M$ is finitely generated. For this we  adapt the proof of Gordan's Lemma \index{Gordan's Lemma} from \cite{Ful} Chap 1. Prop. 1.
 The  set $K=\bcu_{m>0} \frac{1}{m}F_m \cap M $ is a discrete subset of the compact Hausdorff space $F_1$, so is finite. 
Therefore for some $N$, we have $NK\subseteq (F_N \cap M )$.  For example we can take $N$ to be the least common multiple of a finite set of denominators that are required to obtain $K$. Note that $NcK \subseteq M$ for all $c\ge 0.$
We claim that $F\cap M $ is generated by the finite dimensional subspace $\sum_{m \ge 0}^N F_m\cap M .$
Indeed, 
if $v\in F_a\cap M$, then $v/a = s \in K$. Write $a =N c +r$ where $c\ge0$ and  $0\le r < N$. Suppose $v/a = s \in K\subset F_1.$ Then
$v = N cs  +rs $. Since $v, N cs \in M$, it follows that $rs\in F_r\cap M$. This proves the claim.\\ \\
The map $\phi$ extends to a map of free abelian groups 
$\phi:\Z S\lra \Z(F\cap M)$. 
Let $B=\Ker \phi %\{s\in S| \phi(s) =0\}.
$ 
This is a free abelian group which we identify $B$ with $\Z^N$ for some $N$. Then for $I\subseteq [N]$, let $$B_I=
\{(b_1 \ldots, b_N)| b_i \ge 0 \mbox{ iff } i\in I\}.$$ Then $B=B_I= \bcu B_I$.
Each $B_I$ is a finitely generated semigroup. % (Gordan's Lemma again).
Now let $s_1, \ldots ,s_M$ be a finite set of elements of $S$ such that 
$F\cap M$ is generated by the image of this set and $\{s_1, \ldots ,s_M\}$ contains a set of generators for each $B_I$. We claim $s_1, \ldots ,s_M$ generates $S$. If $x\in S$, then $\phi(x) = \sum_i r_i \phi(s_i)$ where $r_i$ is a non-negative integer. Since  $x-  \sum_i r_i s_i \in B_I$ for some $I$ the result follows easily.  
\epf
 \section{Geometry of Šapovalov elements}  \label{GF} 
We have   been doing algebraic geometry all along.
 \ff{ll y a plus de quarante ans que je dis de la prose sans que j'en susse rien, Molière, Le Bourgeois Gentilhomme}
\\ \\
\noi Here we work over $\ttk$ since some of the results we need have a Noetherian hypothesis. %\subsubsection{Geom 1}
\subsection{Sheaves and morphisms}  \label{GF1}
%$\span\\{x^\ttk_{\gs}\}|{\gs \in {{\bf P}}(m\gb)}\}$ ${\pi \in {{\bf P}}(m\gamma)}  e_{-\pi}c_\pi{\gs \in {{\bf P}}(m\gb)}$
Let $f:X\lra Y$ be a morphism of schemes and $\cF, \cG$ quasi-coherent sheaves on $X, Y$. The {\it direct image } \index{Direct image sheaf} $f_* \cF$ of $\cF$ and the {\it inverse image }  \index{Inverse image sheaf}
$f^* \cG$ of $\cG$ are defined in \cite{Ha} page 109. We have a natural isomorphism 
 $$\Hom\cO_X(f^* \cG, \cF) 
\cong  \Hom\cO_Y(\cG, f_*\cF).$$
This is the analog for sheaves of \eqref{GF11}.  

\subsection{A theorem of Serre}  


 Given an $\N$ graded ring $R=\bop_{n\ge 0} R(n)$ with $R(0)$ is a finitely  generated  $\ttk$-algebra,  such that $R$ is generated over  $R(0)$ by $R(1)$,
 we can construct a scheme $X=\pro R$.  The construction gives invertible sheaves $\cO_X(n)$ \cite{Ha} II.2 and II.5.   
%The proof of \cite{Ha} Proposition II.5.12 shows that if $R$ is  finitely generated over  $\ttk$, then the sheaves $\cO_X(n)$  are invertible.  
Denote the category of quasi-coherent sheaves on $X$ by $\frak{Qco}X$.
\index{$\frak{Qco}X$, Category of quasi-coherent sheaves on $X$}     Then by a Theorem of Serre \cite{Se}, $\frak{Qco}X$ is equivalent to the quotient category ${\bf gr} R/\tors$. 
\\\\
In Proposition \ref{td5}  we introduced %We also have the 
graded rings 
$R(\gt,\gs)$ % from Subsection \ref{ta3}
 and in \eqref{zra1}, \eqref{zra2} various isomorphisms between them. Set $\X_{\gt,\gs}= \pro R(\gt,\gs)$ \index{Algebras $R(\gt,\gs)$!$\X_{\gt,\gs}= \pro R(\gt,\gs)$}
and $\X_{\gt}= \pro\mbox{Š}^0_{\gt}$

\bp
\bi \itema We have a commutative diagram of projective schemes

\xymatrix{
&&&&\X_{\gb, \gc}
 \ar@{->}[dr]_{
(B_S\ci\nabla_S^t\ot 1)
}  \ar@{->}[rr]^{
(1\ot P_\ga)} && \X_{\gb, \gb}
\ar@{->}^{f%(\gl, \alpha^\vee)
}[dl] &\\
&&&&&\X_{\gc, \gc}
} 
\itemb
\noi We have equivalences of categories and a commutative diagram. %IAN On the nose?
$$
\xymatrix{\Qcc  \ar@/_1pc/[rr]_{{(B_S\ci\nabla_S^t\ot 1)_{*}}} %\ar@/^/[drr]^x
\ar@{>}^{f_{*}}[r]
& \Qbb \ar@{>}^{{(1\ot P_\ga)}_{*}}[r]&\Qbc}.$$ \ei
\ep\bpf (a) follows since the isomorphisms preserve the irrelevant ideals and (b) holds by 
 Corollary  
\ref{3W2} and the Theorem  of Serre.\epf
\noi By zooming in the map $f$ from (a) induces an isomorphism of fiber products over $\A^1.$
$$\X_\gb \ti \A^1\lra \X_\gc \ti \A^1.$$
Let $\cO_{\gt,\gs}(m) =\cO_{\X_{\gt,\gs}}(m).$ Then sections of 
$\Obb(m)$ have the form
$\sum_{\gs \in {{\bf P}}(m\gb)}x_{\gs} \gk(a_\gs).$
\\ \\Let $Z = (\fn^-)^*$, so that $S(\fn^-)=\cO(Z) $ and $S(\fn^-)_e=\cO(Z)_e = \cO(Z_e)$ where 

\be\label{GF7} 
Z_e = \{z\in Z|e(z)\neq 0\}, \quad V_e= Z \bsk Z_e
\ee
%Next 
%To understand the prime ideals in $\pro  S$, 
Note that 
$\Obb(m)$ and  $(\nabla_S)_*\Obb(m)$  are 
objects in $\Qbb$ and  $\Ocb(m)$ 
is an object in 
$\Qcb$ and we have maps 
\be\label{318}((B_S\ci\nabla_S^t)\ot 1)_*:
 \Obb(m){\lra}  \Obb(m) \lra \Ocb(m)\ee
\be\label{319}(1\ot P_\ga)_*:
\Occ(m)  {\lra}  \Ocb(m)\ee
 Now  $$\gk( \theta_{\gamma,m})=\sum_{\pi \in {{\bf P}}(m\gamma)} 
x_{\pi}\ot \gk( c_\pi) %\in \cM \ot\cO_\gc 
\mbox { and }  
\gk( \theta_{\gb,m})=\sum_{\gs \in {{\bf P}}(m\gb)} x_{\gs}\ot \gk(a_\gs) %\in \cL \ot\cO_\gb
$$
are sections of $\Occ(m), \Obb(m)$ respectively. 
If we evaluate \eqref{318}
and \eqref{319} on $ 
\gk( \theta_{\gb,m})$ and $\gk(\theta_{\gamma,m})$ respectively,
we find exactly the same relation \eqref{1789} as in
Theorem  \ref{zra}, except that algebra maps have been replaced by
direct image functors.
%\vspace{0.2cm}
\subsection{The cancellation step}%egp e7
By \eqref{2W9}, there is a map $F:R(\gb,\gb)\lra S(\fn^-)_e \ot \cO(H_{\gc,0}). 
$ The cancellation step in Lemma  \ref{pre1}  shows $F(\gk(\theta_{\gb, m}))(\gl)\in S(\fn^-)$ for all   
$\gl\in H_{\gc,m}.$  
%Thus  $\theta_{\gamma, m}(\gl)$ has no poles  on $V_e$. 
The proof depends on representation theory.   Let $X =\Spec S(\fn^-)_e \ot \cO(H_{\gc,0})$
and 
$Y=\Spec R(\gb,\gb)$.  The comorphism $F_\#:X\lra Y$ satisfies
$(F \gth)(x) = \gth(F_\#(x))$ for all $\gth\in\cO(Y)$ and all closed points $x\in X.$  In particular for $\gth \in R(\gb,\gb)(m)$ and  $\gl  \in H_{\gc,m}$ $$(F \gth)(\gl) = \gth(F_\#(\gl)).$$ 
The statement $F(\theta)(\gl)\in S(\fn^-)$ can be interpreted in two ways by viewing $\gth$ or $\gl$ as the independent variable.
\noi %Given  $\gth=\sum_{\gs \in {{\bf P}}(m\gb)} x_{\gs}\ot \gk(a_\gs) \in \mbox{Š}_{\gb}(m)\ot\cO(H_{\gb,m})$. 

%\subsubsection{Projection on the first factor}%egp e7

\subsubsection{First interpretation}%egp e7
\noi   
\bp 
Given  $\gth \in \mbox{Š}^0_{\gb}(m)\ot \cO(H_{\gb,m}).$ 
\bi \itema The set
$$K_\gth=\{\gl\in  H_{\gc,m}%\fh^*
|(F\gth)(\gl)\in S(\fn^-)\}$$
is closed.
\itemb If $\gth = \gth_{\gb,m}$, then 
$K_\gth=H_{\gc,m}.$
\ei  \ep \bpf (a) holds since 
$$K_\gth=\{  \gl\in  H_{\gc,m}|\gth (F_\#(\gl))\in V_e
\}.$$ Statement (b) follows from the cancellation step.
\epf 
\subsubsection{Second interpretation}%egp e7
Recall the map 
$y_m^*:H_{\gs, 0} \lra H_{\gs, m}$ from \eqref{jjj}
\bp 
Given $\gl  \in H_{\gc,0}$ % then  $\gth$ a section oF $\Obb(m)$,
\bi \itema There is a graded subalgebra 
$A_\gl= \bop_{m\ge 0}A_\gl(m)$ of $\mbox{Š}^0_{\gb}$ with
$$A_\gl(m)\ot\cO(H_{\gb,m})=\{\gth\in \Obb(m) |(F\gth)(y_m^*\gl)\in S(\fn^-) \}.$$
\itemb 
$A_\gl(m)$ contains $\gth_{\gb,m}$.
\ei\ep \bpf Use Lemma \ref{td5} for (a) and the cancellation step for (b). \epf 

\br {\rm  In type A, 
$K_\gth = \fh^*$ for all $\gth \in \mbox{Š}^0_{\gb}(m)\ot\cO(H_{\gb,m}) $ and 
$A_\gl=\mbox{Š}^0_{\gb}$ for all
$\gl  \in H_{\gc,m}$. 
}\er \subsection{Existence and uniqueness of \v Sapovalov elements.} \label{inc}
\bl Given $\gth\in \Obb(m)$ and $S \subseteq \Gd^+$
$$C_{\gth, S} =\{\mu\in \fh^* |e_\ga \gth v_\mu =0 \mbox{ for all } \ga\in S\}$$
 is closed in $\fh^*$.
\el \bpf This holds by the proof of \cite{D} Lemma 7.6.12.\epf
\noi If $S = \Gd^+$  and $\theta=\theta_{\gb,m}$, then $C_{\gth, S}= {H}_{\gb, m}$\bl 
 Given $S$ and $C$ is closed in $\fh^* $
define $$W_{S, C} =\{\gth\in \Obb(m)|e_\ga \gth v_\mu =0 \mbox{ for all } \ga\in S\}$$
%Each puts some linear conditions on 
If $S= \Gd^+$ and $C = {H}_{\gb, m}$  then $W_{S, C} $ is one dimensional and spanned by $\theta_{\gb,m}$.\el  
\bpf This holds by the existence and uniqueness of \v Sapovalov elements.
\epf
\section{{Examples}}  \label{vf}  
\subsection{Type A}  \label{Type A}  

\bexa \label{PA.13}{\rm
This is the smallest example of a map between \v Sapovalov  algebras $$R=\mbox{Š}_\gb \lra  \mbox{Š}_\gc =S.$$  It arises in the computation of the \v Sapovalov elements for the highest root $\gc$ of $\fg=\fsl(3)$.  
Let $\fg' \cong \fsl(2)$ be the subalgebra of $\fg$ consisting of $2\ti 2$ matrices located in the lower right corner of $\fg.$   Let $\gb$ be the simple root of $\fg'$ and $\ga$ the other simple root of $\fg$. Consider the negative root vectors 
$ X_\ga =Y, X_\gb =X, Z=[X,Y]=X_\gc. $
Then $X,Y, Z$ form a basis for the subalgebra $\fn^-$ of $\fg$.
\\  \\ The \v Sapovalov subalgebras are 
 $R =\ttk[X, ZY^{-1}], S=\ttk[XY, Z],$ with 
$$\deg_R X =\deg_R ZY^{-1}   =1,  \deg_S XY =\deg_S Z = 1 \mbox{ and } \deg T =0.$$ 
We have $d_- = Y^{-1}\pd_Y$, and 
$\nabla_S = \exp(Td_-)=1  +Td_- + \mbox{ terms with higher powers of }T $. Thus  $\nabla_S$ is the automorphism of $R[T]$ given by 
$$ X\lra X +TZY^{-1}, \quad ZY^{-1}\lra ZY^{-1}.$$
\noi
The composite  $B \ci \nabla_S: R[T]\lra S[T]$ satisfies 
$$ X\lra XY +TZ, \quad ZY^{-1}\lra Z.$$
%In the notation of \eqref{121d}, the   \v Sapovalov element$\theta_{\gamma,m}$is given by $(XY +tZ)^m$ modulo terms of lower degree
}\eexa 

\subsection{\v Sapovalov elements for $\fsp(4)$} \label{g19}
Let $\fg =\fsp(4)$ with simple roots $\gb,  \ga+\gb, 2\ga +\gb, \ga$.  Denote the corresponding negative root vectors by $e_{-\gb}, e_{-\ga-\gb}, e_{-2\ga -\gb}, e_{-\ga}$. 
\subsubsection{Case 1}  \label{71a} We find the element $\gth_{\ga+\gb,1}$ for $s_\ga(\gb)= \ga+\gb.$ 
For this we invert $e_{-\ga}$ and then look at the $\gb$ weight space. 
There are two basis elements $e_{-\gb}, e_{-\ga-\gb}e_{-\ga}^{-1}$. We compute that the element
$$\gth_{\ga+\gb,1} = e_{-\ga}e_{-\gb} + (\mbox{ coeff of deg 1})e_{-\ga-\gb}.$$ Thus no rogue partitions occur in $\gth_{\ga+\gb,2} $.
But on the other hand we have an isomorphism of graded algebras $ \mbox{Š}_\gb\lra  \mbox{Š}_{\ga+\gb}$ and $e_{-2\ga -\gb}e_{-\gb}\in \mbox{Š}_{\ga+\gb}(2)$ corresponds to a   rogue partition of $2(\ga+\gb)$.  The explanation for this is that  $ e_{-2\ga -\gb}e_{-\gb}e_{-\ga}^{-2
}$ is an element in degree two in which $B_S$ maps to $e_{-2\ga -\gb}e_{-\gb}$.
\subsubsection{Case 2}  \label{71b}
To find the element $\gth_{2\ga+\gb,1}$ for $s_\gb(\ga)= 2\ga+\gb,$we invert $e_{-\gb}$ and look at the $\ga$ weight space.
There are three basis elements $e_{-\ga}, e_{-\ga-\gb}e_{-\gb}^{-1}, e_{2\ga +\gb}e_{-\gb}^{-2}$. We compute that the element
$$\gth_{2\ga+\gb,1} = e_{-\ga}e_{-\gb}^2 + (\mbox{ coeff of deg 1})e_{-\ga-\gb}e_{-\gb} +(\mbox{ coeff of deg 2})e_{-2\ga -\gb}.$$ 
\subsection{Type C: The cancellation step}  \label{Type C}
%\bexa{\bf An (ortho)-symplectic example} \label{1cosp}\\ \\
%\subsection{Introduction}  \label{xy.1}
%\bexa \label{1eCdeg estimate.4}
{\rm 
\noi A  crucial step in the construction of \v Sapovalov elements is the observation
that the terms 
$ q^{(j)} _{\gs}( e_{-\gz}^\tth)$ in Lemma 
\ref{109b}  
vanish unless %by Lemma \ref{1uebasis},
\begin{eqnarray} \label{198g}
\gz(\ga)\ge N- mq.\nn
\end{eqnarray}
 This follows from representation theory. The condition is necessary since otherwise the evaluation of the \v Sapovalov element would belong to 
$U(\fn^-)_e$ and not to  $U(\fn^-)$.
We give an example where the individual terms $R _{\gs, j}$ on the right of Equation (\ref{xxx}) are not identically zero, and verify directly that  the sum itself is zero, see \eqref{vanish}. %This cannot happen in Type A.  The key difference in the examples below seems to be that it is necessary to apply Equation (\ref{121nd}) more than once with the same simple root $\ga$. 
Consider the Dynkin-Kac diagram below for the Lie superalgebra $\fg = \osp(2,4)$.

\vspace{0.4cm}

\vspace{0.4cm}
\begingroup
\setlength{\unitlength}{0.10in}
\begin{picture}(-30,-10)
\thicklines
\put(14.414,0.0){\line(1,0){9.23}}
\put(25.95,1.0){\line(1,0){10.1}}
\put(25.95,-1.0){\line(1,0){10.1}}
\put(30,0){\line(1,-1){1.5}}
\put(30,0){\line(1,1){1.5}}
\put(12,1){\line(1,-1){2}}%
\put(12,-1){\line(1,1){2}}
\put(13,0){\circle{2.828}}
\put(25,0){\circle{2.828}}
\put(37,0){\circle{2.828}}
\put(11.6,-3.0){$\gep-\gd_1$}
\put(23.1,-3.0){$\gd_1-\gd_2$}
\put(36.2,-3.0){$2\gd_2$}
\end{picture}

\vspace{1.4cm}

\endgroup
\noi If we replace the grey node in the diagram  labelled by the root $\gep-\gd_1$ by a white node, labelled by the root $\ga_0 =\gd_0-\gd_1$ we obtain the Dynkin diagram and root system of $\fsp(6)$.
Note that the equations below involving $\gr$,  \eqref{k1} - \eqref{k3} do not depend on 
$\gep-\gd_1$ or $\ga_0 =\gd_0-\gd_1$. Although the dimensions of the nilradicals 
have dimensions 8 and 9 for $\fg = \osp(2,4)$ and   $\fsp(6)$ respectively, we do not need the root $2\gd_0$ in any  of the computations below.
Another point to make is that  for the case $\fg = \osp(2,4)$ in any partition $\pi$ that we use, there is at most one odd root  $\gc$ such that $\pi(\gc)> 0$. Therefore the commutation relations needed for 
 $\fg = \osp(2,4)$ are exactly the same as  $\fsp(6)$.\\ \\
\noi Let $\fb$ be the Borel subalgebra having  simple
roots  $\gb = \epsilon - \gd_1, \ga_1 = \gd_1 - \gd_2,$ $  \ga_2 = 2\gd_2$ and 
$\fb^-$  the Borel opposite to $\fb$. In the case of 
$\fsp(6)$ set  $\gb = \gd_0- \gd_1.$ 
 Let
$e_{-\gb}, e_{-\ga_1}, e_{-\ga_2}  $ be the negative simple root vectors.
We compute the \v Sapovalov elements
$\gth_1, \gth_2, \gth_3$ for the positive roots $\gc_1= \gb + \ga_1, \gc_2= \gb + \ga_1+ \ga_2$ and $\gc_3= \gb + 2\ga_1+ \ga_2$.
%A remarkable fact is that all the coefficients of these\v Sapovalov elements are products of linear factors.  Furthermore the vanishing of these linear factors corresponds to factorizations of the \v Sapovalov elements.% as in \cite{M21}.
\subsubsection{\v Sapovalov elements}  \label{xy.2}
%\pagebreak
Define the other negative root vectors by
$$e_{- \ga_1 - \ga_2} = [e_{-\ga_1},e_{- \ga_2 }  ],  \quad \quad
e_{- 2\ga_1 - \ga_2} = [e_{-\ga_1}, e_{- \ga_1 - \ga_2}],  $$
$$e_{-\gb- \ga_1} =[e_{-\ga_1} ,e_{-\gb}
], \quad e_{-\gb- \ga_1-\ga_2} = [e_{- \ga_2 } ,e_{-\gb- \ga_1}], \quad e_{-\gb- 2\ga_1 - \ga_2} = [e_{-\ga_1},e_{-\gb- \ga_1-\ga_2}]. $$
It follows from the Jacobi identity that
$$[e_{-\gb}, e_{- \ga_1 - \ga_2}] = e_{-\gb- \ga_1-\ga_2},
\quad \quad
[e_{- \ga_1 - \ga_2},e_{-\gb- \ga_1}] = e_{-\gb- 2\ga_1 - \ga_2},$$ and
$$[e_{-\gb},e_{- 2\ga_1 - \ga_2}] = 2e_{-\gb- 2\ga_1-\ga_2}.$$
\noi 
Throughout we order the set of positive roots so that 
if $\gs$ belongs to the set
\be\label{X1}\{
\gb, \gb + \ga_1, \gb + \ga_1 + \ga_2, \gb + 2\ga_1 + \ga_2\}\ee 
%Next order the set of positive roots so that 
if $\pi$ is a partition   
with $\pi(\gs)> 0$ then  $e_{-\gs}$ is the 
first factor in $e_{-\pi}$.
\\ \\
Let $s_1, s_2$ be the reflections corresponding to the simple roots $\ga_1, \ga_2.$
Then for $\gl \in \fh^*$ define $\gl_1 =s_1\cdot\gl, \;\gl_2 = s_2\cdot\gl_1, \;\gl_3 = s_1\cdot\gl_2$. Using the notation \eqref{efh}, set $t_i= t_{\ga_i}$ for $i=1,2$. For a root $\gc$ any polynomial function on ${H}_{\gc}$  can be expressed uniquely as a polynomial in $t_1, t_2.$
Let
\be\label{k1}p=(\gl + \gr,\ga^\vee_1) = 
-(\gl_1 + \gr,\ga^\vee_1) =  -(\gl_2 + \gr,(\ga_1+\ga_2)^\vee)=-(\gl_3 + \gr,(\ga_1+\ga_2)^\vee),\ee 
\be\label{k2}
q= (\gl + \gr,(2\ga_1  + \ga_2)^\vee) = (\gl_1 + \gr,\ga^\vee_2) = -(\gl_2 + \gr,\ga_2^\vee) = -(\gl_3 + \gr,(2\ga_1+\ga_2)^\vee) ,\ee 
and 
\be\label{k3}r=(\gl + \gr,(\ga_1+\ga_2)^\vee) =(\gl_1 + \gr,(\ga_1+\ga_2)^\vee)  =
(\gl_2 + \gr,\ga^\vee_1)  = -(\gl_3 + \gr,\ga_1^\vee) .\ee
Then $r = 2q - p.$  
%Let $\gc_1 =  \gb + \ga_1,  \gc_2 =  \gb + \ga_1 +\ga_2$ and $\gc_3 =  \gb + 2\ga_1 +\ga_2.$ 
We compute the
\v Sapovalov elements
$\gth_{\gc_i}$ for $\fg$. 
To do this we use Equation (\ref{121nd}) and  the abbreviations $\gth_{i}: = \gth_{\gc_i}(\gl_i)$.   %We can assume $\gc \neq \gb.$ 
To start the induction, assume that $(\gl+\gr,\gb) =(\gb,\gb)/2.$ Suppose that $p, q, r$ are nonnegative integers. Then
\be \label{nelab}
e_{-\ga_1}^{p+1}e_{-\gb} = \gth_1e_{-\ga_1}^{p},\quad e_{- \ga_2 }^{q+1}\gth_1 = \gth_2e_{- \ga_2 }^{q},\quad 
e_{- \ga_1 }^{r+1}\gth_2 = \gth_3e_{- \ga_1 }^{r}.\ee 
In the Verma module $M(\gl)$ induced from $U(\fb)$ with highest weight $\gl$ and highest weight vector $v_\gl$, set $v_1=e_{-\ga_1}^{p}v_\gl,$ $v_{ _2  }=e_{- \ga_2 }^{q}v_1$ and $v_{3}= e_{- \ga_1 }^{r}v_2$. These are all highest weight vectors for 
$\fb$ and by \cite{M101}, Theorem 9.3.2 they generate Verma submodules of $M(\gl)$. 
In $U(\fb^-)$ we have 
$$[e_{-\ga_1}^{p+1} ,e_{-\gb}
] = (p+1)e_{-\gb- \ga_1} e^p_{-\ga_1}$$
$$[e_{- \ga_2 }^{q+1} ,e_{-\gb- \ga_1}] =
(q+1)e_{-\gb- \ga_1-\ga_2} e_{- \ga_2 }^q$$
$$[e_{- \ga_2 }^{q+1} ,e_{-\ga_1}]  =-(q+1)e_{- \ga_1 - \ga_2}
e_{- \ga_2 }^q.$$
This easily gives
\be \label{1tha}
\gth_1 = (p+1)e_{-\gb- \ga_1} + e_{-\gb}e_{-\ga_1} = pe_{-\gb- \ga_1} + e_{-\ga_1}e_{-\gb}.\ee
We order the set of positive roots so that for any  partition $\pi$, $e_{-\ga_2}$ occurs last if at all in $e_{-\pi}$. Then %
\begin{eqnarray}\label{1la2}
\gth_2
&=& (p+1)[(q+1)e_{- \gb - \ga_1 - \ga_2} +e_{-\gb- \ga_1}e_{-\ga_2}] +e_{-\gb}[ e_{-\ga_1}e_{- \ga_2} - (q+1)e_{- \ga_1 - \ga_2} ].\nn\\
&=& (p+1)[(q+1)e_{- \gb - \ga_1 - \ga_2} +e_{-\gb- \ga_1}e_{-\ga_2}] +e_{-\gb}[ e_{-\ga_2}e_{- \ga_1} - qe_{- \ga_1 - \ga_2} ].
\end{eqnarray}
Next order the set of positive roots so that for any  partition $\pi$, $e_{-\ga_1}$ occurs last $e_{-\pi}.$ To find $\gth_3$  we use
$$[e_{-\ga_1}^{r+1} ,e_{-\gb-\ga_1-\ga_2}
] = (r+1)e_{-\gb- 2\ga_1 -\ga_2} e^r_{-\ga_1},$$
$$[e_{-\ga_1}^{r+1} ,e_{-\gb-\ga_1}e_{-\ga_2}
] = (r+1)e_{-\gb- \ga_1}e_{-\ga_1-\ga_2}e^
{r}_{-\ga_1}
+\left( \begin{array}{c}
                r+1 \\
                2 \end{array}\right)e_{-\gb- \ga_1}e_{-2\ga_1 -\ga_2} e^{r-1}_{-\ga_1},$$
$$[e_{-\ga_1}^{r+1} ,e_{-\gb}e_{-\ga_1-\ga_2}
] = (r+1)[e_{-\gb}e_{- 2\ga_1 -\ga_2} e^r_{-\ga_1}
 +e_{-\gb- \ga_1}e_{-\ga_1 -\ga_2} e^r_{-\ga_1} +re_{-\gb-\ga_1}e_{-2\ga_1-\ga_2}
e_{\ga_1}^{r-1}],$$
$$e_{-\ga_1}^{r+1} e_{-\gb}e_{-\ga_2}e_{-\ga_1}
 =
e_{-\gb}[e_{-\ga_2}
e_{- \ga_1}^{2} +
(r+1)e_{- \ga_1-\ga_2}e_{-\ga_1}
+\left( \begin{array}{c}
                r+1 \\
                2 \end{array}\right)e_{- 2\ga_1 -\ga_2}] e^{r}_{-\ga_1}$$
$$+ (r+1)e_{-\gb- \ga_1}[e_{-\ga_2}
e_{- \ga_1} +
re_{- \ga_1-\ga_2}]e^{r}_{-\ga_1}
+(r-1)\left( \begin{array}{c}
                r+1 \\
                2 \end{array}\right)e_{-\gb-\ga_1}e_{-2\ga_1-\ga_2}
e_{\ga_1}^{r-1}
                %e_{-\gb}e_{- 2\ga_1 -\ga_2} e^{r-1}_{-\ga_1}
                .$$
The above equations allow us to write $e_{- \ga_1 }^{r+1}\gth_2$  in terms of elements $e_{-\pi}$ with $\pi$ a partition of $\gb + (r+2)\ga_1 + \ga_2$.
We see that the term $e_{-\gb- \ga_1}e_{- 2\ga_1 -\ga_2} e^{r-1}_{-\ga_1}$  occurs in $e_{- \ga_1 }^{r+1}\gth_2$ with coefficient
\be \label{vanish}
\left( \begin{array}{c}
                r+1 \\
                2 \end{array}\right)[(p+1) -2q + (r-1)] = 0.\ee
This is predicted by the cancellation step in Subsection \ref{pre}. %the proof of Lemma \ref{1wpfg}. 
In the remaining terms, $e_{- \ga_1 }^{r}$ can be factored on the right, and this yields
\begin{eqnarray}\label{1th3}
\gth_3 &=&
(p+1)(q+1)(r+1)e_{-\gb- 2\ga_1 -\ga_2} +(p+1)(q+1)e_{-\gb- \ga_1-\ga_2}e_{-\ga_1} \\
&+& (q+1)(r+1)e_{-\gb- \ga_1}e_{-\ga_1 -\ga_2}
%\\&& might be q(r+1)e_{-\gb- \ga_1}e_{-\ga_1 -\ga_2}
-(p/2)(r+1)e_{-\gb}e_{- 2\ga_1 -\ga_2}\nonumber \\
&+& 2(q+1)e_{-\gb- \ga_1}e_{-\ga_2}e_{-\ga_1}+ (r-q+1)e_{-\gb}e_{-\ga_1-\ga_2}e_{-\ga_1}+e_{-\gb}e_{-\ga_2}
e_{- \ga_1}^{2}. \nonumber
\end{eqnarray}


\subsection{Type C: Decomposition into simple $\fl$-modules}  \label{Type C2}
Given the importance of the adjoint action of $\fl$, we might hope 
to decompose each  graded component  $\mbox{Š}_{\gb,\Z}(m)$  is a direct sum of simple $\fl$-modules (over $\Z$). 
We focus on the $\fsl(2)$ case and indicate how this can be done, see Example \ref{nB2} for the case of $\fsp(6)$ when $mq=1$. However the process depends on the choice of a subset of 
${{\bf P}}(m\gb)$ and there can also be highest weight vectors which have poles, see Example \ref{nB2} for the case of $\fsp(6)$ when $mq=1$.
Hence outside of Type A, it seems unlikely that there is a canonical way to do this.

\subsubsection{How to improve a $\Z $-basis}  
In Subsection \ref{cs} we gave 	a 
$\Z $-basis for the weight space $S_\Z(\fn^-)_e^{-\eta}.$
 and hence for each graded component  
$\mbox{Š}_{\gb,\Z}(m)$. Let $\fsl(2)$ be the  Lie algebra with basis $E, F, H$ such that $[H,E]=2E, [H,F]=-2F$ and  $[E,F]=H$. Denote the simple $\fsl(2)$ module of dimension $n+1$ by  $L(n)$. In the adjoint representation, $E, F$ act as $d_{+}, d_{-}$ respectively and $H$ acts as $\ad h_{\ga}.$  %In terms of the notation from the introduction to Section \ref{PA.2}, we have \be \label{E7} F = d_-= x_\ga^{-1} \ad e_{-\ga} = x_\ga^{- 1} \pd_S.\ee
We want to improve the $\Z $-basis  given by Lemma \ref{1uebasis} 

\bt \label{Z-b} Let  \be M=\bigoplus_{n\in \Z} M(n)\ee  be a $\Z$ lattice in a finite dimensional $\fsl(2)$-module and $u\in M(n)$. 
Assume that $M$ is stable under $E^{(k)}, F^{(k)}$ for all $k$. 
Then there is a unique expression 
\be \label{arl}u = \sum_{k\ge 0} F^{(k)} u(k)\ee where $u(k) \in M(n+2k)\cap \ker E.$
\et 
\bpf This is the limit as $q\lra 1$ of \cite{HK} Lemma 4.1.1.  It is not hard to give a direct proof.
\epf 
\noi The Theorem  shows  
any weight vector $v$ of positive weight is congruent to a highest weight vector  $v_0$ mod the image of $F$. 
\noi For $\cB \subseteq M$, let $\langle \cB \rangle$ be the $\fsl(2)$ submodule generated by $\cB$-
\bc \label{ZB}
 Suppose $\cB= \ds \cB_k$  be a $\Z$-basis for $M(n+2k)\cap \ker E$ for all $k>0$ and let $\cC_0$ be a $\Z$-basis for $M(n)$. % $\Z$ lattice in a finite dimensional $\fsl(2)$-module and $u\in M(n)$. Then there is a unique expression
Then there is a $\Z$-basis $\cB_0$ of $M(n)$ consisting of  highest weight vectors such that 
$\langle \cB_0 \cup \cB \rangle= \langle \cC_0 \cup \cB \rangle$
\ec \bpf This follows directly from Theorem \ref{Z-b}.
  \epf \noi This allows us to write  $M$ as a direct sum of simple modules.  We start by collecting the modules of largest dimension and then proceed to lower dimension.  The process is probably best understood from an example. 
\bexa \label{nB2} {\rm 
Let  $\fg =\fsp(6)$ with simple roots $\ga_0, \ga_1, \ga_2$ where $\ga_2$ is a long root and $\ga_0+\ga_2$ is not a root.   Let $x_0, x_1, x_2$ be root vectors with $\wt\, x_i = -\ga_i$.  If $\gb=\ga_0+ \ga_1+ \ga_2$ and $\gc=\ga_0+ 2\ga_1+ \ga_2$, we decompose $M=$Š$_\gb(1)$ into a direct sum of simple modules. This example is essentially the same as in Subsection \ref{Type C}  where $\gb, \gc$ were called $\gc_2$ and $\gc_3$. 
 We have $mq=1.$
 There are 4 partitions $\gs$ of $\gb$ and we order the corresponding $x_\gs$ as    
\be \label{1bc} x_0 x_1x_2,\quad [x_0, x_1]x_2,\quad x_0 [x_1,x_2],\quad [x_0, [x_1,x_2]].\ee %\def\rer{{\operatorname{res}}}
 If $x_\gs$ is the first of these, then $x_\gs$
 is a highest weight vector of weight 3 since $d_+(e_{-\ga}) =0$.
 From this   
we obtain the summand $\langle x_0 x_1x_2 \rangle$
 which is isomorphic to $L(3)$.  The last 
 $w=[x_0 [x_1,x_2]]$  
is a highest  vector so  $w_0=w$ and %=[x_0 [x_1,x_2]]$
$M$ contains $\langle  w_0 \rangle.$ 
Consider the second and third terms  $u=[x_0, x_1]x_2,\; v=x_0 [x_1,x_2]$.   In the notation of \eqref{arl}, $u_0=\pm v_0$ 
and we need to choose one of $u_0, v_0$ to include in the sum.  Let us choose $u_0$. Then 
\by \label{107bc} \langle x_0 x_1x_2 \rangle  \op \langle u_0  
\rangle \op \langle w_0  \rangle \subset M. \ey 
We show that $M$ contains one more submodule of dimension 2. This is most easily done by looking for lowest weight vectors.
Now $(S(\fn^-)_e)^{\gb}$ has a basis consisting of all products 
$e_{-\pi} e^k$ 
where  $k \in \mathbb{Z}$ and $\pi \in {\overline{\bf P}}(\gb -k\ga)$ with  $\pi(\ga) = 0.$  It is easily seen that ${\overline{\bf P}}(\gb -k\ga)$ is empty  unless $k\in [-2,1]$.
If $k\ge 0,$ then there are 4 partitions with this property and the corresponding monomials are listed in \eqref{1bc}. It remains to consider partitions with poles. 
 If $k\in [-2,-1]$ we find 4 partitions 
associated to the monomials below, all with weight $-1$.
$$[x_0, x_1][x_1,x_2]e^{-1},\quad x_0 [x_1, [x_1,x_2]]e^{-1},\quad [x_1,[x_2,[x_0, x_1]] ]e^{-1},  [x_0, x_1][x_1 [x_1,x_2]]e^{-2} $$ Since the sum on the left side of \eqref{107bc} has only a 3 dimensional $-1$ weight space, there is a lowest weight vector $z$ of weight $-1$ and we have 
\by \label{107bd} 
M &=& \langle x_0 x_1x_2 \rangle  \op \langle u_0  
\rangle \op \langle w_0  \rangle \op \langle z  \rangle
\nn\\ & \cong & L(3)\op 3L(1).\nn
\ey  }\eexa


\printindex
 
\begin{bibdiv}
\begin{biblist}
\label{bib}

\bib{BGG1}{article}{
%%   author={Bern{\v{s}}te{\u\i}n, I. N.},
   author={Bernstein, I. N.},
   author={Gel{\cprime}fand, I. M.},
   author={Gel{\cprime}fand, S. I.},
   title={Structure of representations that are generated by vectors of
   higher weight},
   language={Russian},
   journal={Funckcional. Anal. i Prilo\v zen.},
   volume={5},
   date={1971},
   number={1},
   pages={1--9},
   issn={0374-1990},
   review={\MR{0291204 (45 \#298)}},
}


\bib{BGG2}{article}{
%%   author={Bern{\v{s}}te{\u\i}n, I. N.},
   author={Bernstein, I. N.},
   author={Gel{\cprime}fand, I. M.},
   author={Gel{\cprime}fand, S. I.},
   title={Differential operators on the base affine space and a study of
   ${\germ g}$-modules},
   conference={
      title={Lie groups and their representations (Proc. Summer School,
      Bolyai J\'anos Math. Soc., Budapest, 1971)},
   },
   book={
      publisher={Halsted, New York},
   },
   date={1975},
   pages={21--64},
   review={\MR{0578996 (58 \#28285)}},
}

		

\bib{BH}{book}{
   author={Bruns, Winfried},
   author={Herzog, J\"{u}rgen},
   title={Cohen-Macaulay rings},
   series={Cambridge Studies in Advanced Mathematics},
   volume={39},
   publisher={Cambridge University Press, Cambridge},
   date={1993},
   pages={xii+403},
   isbn={0-521-41068-1},
   review={\MR{1251956}},
}
		


\bib{D}{book}{
   author={Dixmier, Jacques},
   title={Enveloping algebras},
   series={Graduate Studies in Mathematics},
   volume={11},
   note={Revised reprint of the 1977 translation},
   publisher={American Mathematical Society},
   place={Providence, RI},
   date={1996},
   pages={xx+379},
   isbn={0-8218-0560-6},
   review={\MR{1393197 (97c:17010)}},
}
%\bib{E}{book}{author={Eisenbud, David},   title={Commutative algebra},  series={Graduate Texts in Mathematics},  volume={150},  note={With a view toward algebraic geometry},   publisher={Springer-Verlag, New York},  date={1995},   pages={xvi+785},  isbn={0-387-94268-8},   isbn={0-387-94269-6},   review={\MR{1322960}},  doi={10.1007/978-1-4612-5350-1},}


\bib{FG}{article}{
   author={Fioresi, R.},
   author={Gavarini, F.},
   title={Chevalley supergroups},
   journal={Mem. Amer. Math. Soc.},
   volume={215},
   date={2012},
   number={1014},
   pages={vi+64},
   issn={0065-9266},
   isbn={978-0-8218-5300-9},
   review={\MR{2918543}},
   doi={10.1090/S0065-9266-2011-00633-7},
}


\bib{Ful}{book}{
   author={Fulton, William},
   title={Introduction to toric varieties},
   series={Annals of Mathematics Studies},
   volume={131},
   note={The William H. Roever Lectures in Geometry},
   publisher={Princeton University Press},
   place={Princeton, NJ},
   date={1993},
   pages={xii+157},
   isbn={0-691-00049-2},
   review={\MR{1234037 (94g:14028)}},
}

\bib{G3}{article}{
   author={Gorelik, Maria},
   title={On the ghost centre of Lie superalgebras},
   journal={Ann. Inst. Fourier (Grenoble)},
   volume={50},
   date={2000},
   number={6},
   pages={1745--1764 (2001)},
   issn={0373-0956},
   review={\MR{1817382 (2002c:17017)}},
}


\bib{Gk}{article}{ author={Gorelik, Maria}, title={The Kac construction of the centre of $U(\germ g)$ for Lie superalgebras}, journal={J. Nonlinear Math. Phys.}, volume={11}, date={2004}, number={3}, pages={325--349}, issn={1402-9251}, review={\MR{2084313 (2005f:17011)}}, }


\bib{Ha}{book}{ author={Hartshorne, Robin}, title={Algebraic geometry}, note={Graduate Texts in Mathematics, No. 52}, publisher={Springer-Verlag}, place={New York}, date={1977}, pages={xvi+496}, isbn={0-387-90244-9}, review={\MR{0463157 (57 \#3116)}}, }

\bib{HK}{book}{
   author={Hong, Jin},
   author={Kang, Seok-Jin},
   title={Introduction to quantum groups and crystal bases},
   series={Graduate Studies in Mathematics},
   volume={42},
   publisher={American Mathematical Society, Providence, RI},
   date={2002},
   pages={xviii+307},
   isbn={0-8218-2874-6},
   review={\MR{1881971}},
   doi={10.1090/gsm/042},
}

\bib{H}{book}{ author={Humphreys, James E.}, title={Introduction to Lie algebras and representation theory}, note={Graduate Texts in Mathematics, Vol. 9}, publisher={Springer-Verlag}, place={New York}, date={1972}, pages={xii+169}, review={\MR{0323842 (48 \#2197)}}, }



%\bib{H1}{book}{ author={Humphreys, James E.}, title={Linear algebraic groups}, note={Graduate Texts in Mathematics, No. 21}, publisher={Springer-Verlag}, place={New York}, date={1975}, pages={xiv+247}, review={\MR{0396773 (53 \#633)}}, }


\bib{H2}{book}{
   author={Humphreys, James E.},
   title={Representations of semisimple Lie algebras in the BGG category
   $\scr{O}$},
   series={Graduate Studies in Mathematics},
   volume={94},
   publisher={American Mathematical Society},
   place={Providence, RI},
   date={2008},
   pages={xvi+289},
   isbn={978-0-8218-4678-0},
   review={\MR{2428237}},
}


\bib{J1}{book}{ author={Jantzen, Jens Carsten}, title={Moduln mit einem h\"ochsten Gewicht}, language={German}, series={Lecture Notes in Mathematics}, volume={750}, publisher={Springer}, place={Berlin}, date={1979}, pages={ii+195}, isbn={3-540-09558-6}, review={\MR{552943 (81m:17011)}}, }
\bib{Kac1}{article}{ author={Kac, V. G.}, title={Lie
superalgebras}, journal={Advances in Math.}, volume={26},
date={1977}, number={1}, pages={8--96}, issn={0001-8708},
review={\MR{0486011 (58 \#5803)}}, }



\bib{Kac4}{article}{ author={Kac, V. G.}, title={Laplace operators of infinite-dimensional Lie algebras and theta functions}, journal={Proc. Nat. Acad. Sci. U.S.A.}, volume={81}, date={1984}, number={2, Phys. Sci.}, pages={645--647}, issn={0027-8424}, review={\MR{735060 (85j:17025)}}, }

\bib{KK}{article}{ author={Kac, V. G.}, author={Kazhdan, D. A.}, title={Structure of representations with highest weight of infinite-dimensional Lie algebras}, journal={Adv. in Math.}, volume={34}, date={1979}, number={1}, pages={97--108}, issn={0001-8708}, review={\MR{547842 (81d:17004)}}, }		

\bib{KL}{article}{
   author={Kumar, Shrawan},
   author={Letzter, Gail},
   title={Shapovalov determinant for restricted and quantized restricted
   enveloping algebras},
   journal={Pacific J. Math.},
   volume={179},
   date={1997},
   number={1},
   pages={123--161},
   issn={0030-8730},
   review={\MR{1452529}},
   doi={10.2140/pjm.1997.179.123},
}		


\bib{Mq}{book}{
   author={Marquis, Timoth\'{e}e},
   title={An introduction to Kac-Moody groups over fields},
   series={EMS Textbooks in Mathematics},
   publisher={European Mathematical Society (EMS), Z\"{u}rich},
   date={2018},
   pages={xi+331},
   isbn={978-3-03719-187-3},
   review={\MR{3838421}},
   doi={10.4171/187},
}

\bib{Md}{article}{author={Mudrov, Andrei}, title={Factorization of Shapovalov elements},journal={prepint, arXiv:2202.06220},date={2022},}

\bib{M101}{book}{
   author={Musson, Ian M.},
   title={Lie superalgebras and enveloping algebras},
   series={Graduate Studies in Mathematics},
   volume={131},
   publisher={American Mathematical Society},
   place={Providence, RI},
   date={2012},
   pages={xx+488},
   isbn={978-0-8218-6867-6},
   review={\MR{2906817}},
}


\bib{M21}{article}{author={Musson, I.M.}, title={
\v
Sapovalov elements and the Jantzen sum formula for
contragredient Lie superalgebras}, date={2017},
journal={arXiv:math/1710.10528v1.}}

\bib{M23}{article}{
   author={Musson, Ian M.},
   title={Explicit expressions for \v{S}apovalov elements in Type A},
   journal={J. Algebra},
   volume={623},
   date={2023},
   pages={358--394},
   issn={0021-8693},
   review={\MR{4557791}},
   doi={10.1016/j.jalgebra.2022.12.033},
}
	

\bib{Sh}{article}{ author={{\v{S}}apovalov, N. N.}, title={A certain bilinear form on the universal enveloping algebra of a complex semisimple Lie algebra}, language={Russian}, journal={Funkcional. Anal. i Prilo\v zen.}, volume={6}, date={1972}, number={4}, pages={65--70}, issn={0374-1990}, review={\MR{0320103 (47 \#8644)}}, }


\bib{S1}{article}{author={Serganova, V.},title={Kac-Moody superalgebras and integrability},conference={title={Developments and trends in infinite-dimensional Lie theory},},book={series={Progr. Math.},volume={288},publisher={Birkh\"auser Boston Inc.},place={Boston, MA},},date={2011},pages={169--218},review={\MR{2743764 (2011m:17056)}}}%   doi={10.1007/978-0-8176-4741-4_6},}
		

\bib{Se}{article}{
   author={Serre, Jean-Pierre},
   title={Faisceaux alg\'{e}briques coh\'{e}rents},
   language={French},
   journal={Ann. of Math. (2)},
   volume={61},
   date={1955},
   pages={197--278},
   issn={0003-486X},
   review={\MR{68874}},
   doi={10.2307/1969915},
}
 \end{biblist}
\end{bibdiv}

%\input{2Paper1}




\end{document}
Compare \eqref{q11} with $F=T^m$, $G=T^n$, $ \frac{\ga}{(\ga,\gb^\vee)}=\varpi =q, \mu=p$. 
$$(F\cdot G)(\mu-(m+n)\varpi) =F(\mu-m\varpi)G(\mu-n\varpi)$$
becomes
$$(F\cdot G)(p-(m+n)q) =F(p-mq)G(p-nq)$$
But in Lemma  \ref{11768} we had 
% $\gl \in H_{\gc,m}, (\mu,\gb^\vee)=m, (\gl,\gc^\vee)=m, \;\alpha \in\Pi_{\rm nonisotropic} \cap A(\mu)$ 
 $\mu \in H_{\gb,m}, (\mu,\gb^\vee)=m$ %(\gl,\gc^\vee)=m, \;
%\alpha \in\Pi_{\rm nonisotropic} \cap A(\mu)$ 
Substitute following into \eqref{q11}
\be %\gl = s_\ga\cdot\mu,\;\;\gc = s_\alpha\gb,\;\;
p = (\mu, \alpha^\vee),\;\;q = (\gc, \alpha^\vee)=-(\gb, \alpha^\vee).\ee
Note that
$$(\mu-m\varpi,\ga^\vee) =p + mq$$


\noi In general, let $N(m\gc)$  be the subspace of $\mbox{Š}_\gc(m)$ spanned by all nonrogue partitions of $m\gc$. Then  \be \label{AB} \N_\gc =\bop_{m>0} N(m\gc)\ee is a graded ideal of 
$\mbox{Š}_\gc$.    
There is also  a graded subalgebra $R_\gc$ generated by rogue partitions
and 
$\mbox{Š}_\gc=R_\gc\op \N_\gc$.\\ \\

IAN Projections: Below is the picture for rings



\noi Both inclusions are split because $ \cO(H_{\gs,0})$
is the zero component of  
$R(\gt,\gs)$ and using the notation of \eqref{RA}
$\mbox{Š}^0_{\gt} = \bop_{m\ge 0} \mbox{Š}^0_{\gr} (m)\ot 1$

\xymatrix{
&&&&& R(\gt,\gs)\\
&&&&\mbox{Š}^0_{\gt}
 \ar@{->}[ur]_{
%(B_S\ci\nabla_S^t\ot 1)
}  &&  \cO(H_{\gs,0})\ar@{->}^{%f%(\gl, \alpha^\vee)
}[ul] &
%\\&&&&&\tilde R(\gt,\gs)%\X_{\gc, \gc}
}

	
 $A$ and  $x_1, \ldots, x_m \in A$.
Using divided powers $\pd^{(a)} = \pd^{a}/a!$ we can write the Leibniz rule in the form
\be\label{2ld} \pd^{(n)} (x_1 \ldots x_m ) =\sum  %\left(\begin{array}{c} n\\ a_1, \ldots, a_m \end{array} \right)
\pd^{(a_1)} (x_1) \ldots \pd^{(a_m)}(x_m).
\ee 
 where the sum is over all   $m$-tuples of  nonnegative integers $a_1, \ldots, a_m$ with  $$a_1+ \ldots+ a_m =n.$$


\noi \bl We have a fiber product 

\xymatrix{
&&&&& \X_{\gt,\gs}=\pro R(\gt,\gs)\\
&&&&\pro \mbox{Š}^0_{\gt}
 \ar@{<-}[ur]^{p
%(B_S\ci\nabla_S^t\ot 1)
}  &&  H_{\gs,0}\ar@{<-}_{q%f%(\gl, \alpha^\vee)
}[ul] &
%\\&&&&&\tilde R(\gt,\gs)%\X_{\gc, \gc}
}
\el \bpf	 In general if $A\lra B$ is an inclusion of graded rings such that $B_+ = BA_+,$ then we have a fiber product
$$\pro B =\pro A \ti \Spec B(0).$$
\epf